% =============================================================================
% OPTIMASI BETON BERTULANG PADA STRUKTUR PORTAL RUANG
% Tugas Akhir Sarjana Strata Satu -- Yohan Naftali (07712/TSS), 1999
% Universitas Atma Jaya Yogyakarta
%
% LaTeX transcription of "Optimasi Beton/Teori/*.doc" (legacy MS Word 97-2003
% binary files). See LEGACY_TRANSCRIPTION_NOTES.md in this directory for the
% transcription methodology, what was recovered verbatim, what was
% reconstructed from the accompanying C++ source code, and what could not be
% recovered at all (figures, scanned diagrams, the Tabel 2-1 matrix image).
% =============================================================================
\documentclass[12pt,a4paper]{report}

% ---- Encoding / fonts ----
\usepackage[utf8]{inputenc}
\usepackage[T1]{fontenc}
\usepackage{lmodern}
\usepackage{textcomp}

% ---- Page layout (standard Indonesian thesis margins) ----
\usepackage[a4paper,left=4cm,right=3cm,top=4cm,bottom=3cm]{geometry}
\usepackage{setspace}
\onehalfspacing

% ---- Math / symbols ----
\usepackage{amsmath}
\usepackage{amssymb}
\usepackage{amsfonts}

% ---- Tables / figures ----
\usepackage{graphicx}
\usepackage{array}
\usepackage{longtable}
\usepackage{booktabs}
\usepackage{caption}
\usepackage{multirow}
\usepackage{rotating}
\usepackage{pdflscape}

% ---- Misc ----
\usepackage{enumitem}
\usepackage{titlesec}
\usepackage{fancyhdr}
\usepackage{xcolor}
\usepackage[hidelinks]{hyperref}

% ---- Indonesian captions (no babel "bahasa" dependency assumed) ----
\renewcommand{\contentsname}{DAFTAR ISI}
\renewcommand{\listfigurename}{DAFTAR GAMBAR}
\renewcommand{\listtablename}{DAFTAR TABEL}
\renewcommand{\chaptername}{BAB}
\renewcommand{\figurename}{Gambar}
\renewcommand{\tablename}{Tabel}

% ---- Chapter heading style: "BAB I" / "PENDAHULUAN" centered, like the original ----
\titleformat{\chapter}[display]
  {\centering\normalfont\bfseries}
  {\chaptertitlename\ \Roman{chapter}}
  {6pt}
  {\Large}
\titlespacing*{\chapter}{0pt}{0pt}{30pt}

\titleformat{\section}{\normalfont\bfseries}{\thesection}{1em}{}
\titleformat{\subsection}{\normalfont\bfseries}{\thesubsection}{1em}{}

% ---- Header/footer ----
\pagestyle{fancy}
\fancyhf{}
\fancyfoot[C]{\thepage}
\renewcommand{\headrulewidth}{0pt}

% ---- Placeholder box for figures/tables that were embedded images in the
%      original .doc and could not be recovered from the binary format ----
\newcommand{\gambarhilang}[2]{%
  \begin{center}
  \fbox{\parbox{0.85\textwidth}{\centering\vspace{2cm}
    \textit{[Gambar/citra asli tidak dapat dipulihkan dari berkas \texttt{.doc} --
    sisipkan ulang dari sumber asli]}\\[0.5em]
    \textbf{#1}\vspace{2cm}}}
  \captionof{figure}{#2}
  \end{center}
}

\begin{document}

% =============================================================================
% BAGIAN AWAL (FRONT MATTER) -- halaman diberi nomor angka Romawi kecil
% =============================================================================
\pagenumbering{roman}

% Halaman Judul (title page) -- transcribed from Cover/Cover Depan.doc
\begin{titlepage}
\begin{center}
\vspace*{1.5cm}
{\Large\bfseries OPTIMASI BETON BERTULANG PADA STRUKTUR PORTAL RUANG}

\vspace{1.5cm}
{\large TUGAS AKHIR SARJANA STRATA SATU}

\vspace{3cm}
Oleh :

\vspace{0.5cm}
{\large\bfseries YOHAN NAFTALI}

\vspace{0.3cm}
No. Mahasiswa : 07712/TSS \\
NIRM : 950051053114120051

\vfill

{\bfseries PROGRAM STUDI TEKNIK SIPIL}\\
FAKULTAS TEKNIK\\
UNIVERSITAS ATMA JAYA YOGYAKARTA\\
YOGYAKARTA\\
AGUSTUS 1999

\end{center}
\end{titlepage}

% Halaman Pengesahan (approval page) -- transcribed from Cover/Cover Depan.doc
\thispagestyle{empty}
\begin{center}
{\bfseries PENGESAHAN}

\vspace{1cm}
Tugas Akhir Sarjana Strata Satu

\vspace{0.3cm}
{\bfseries OPTIMASI BETON BERTULANG PADA STRUKTUR PORTAL RUANG}

\vspace{1.5cm}
Oleh :

\vspace{0.3cm}
{\bfseries YOHAN NAFTALI}

No. Mahasiswa : 07712 / TSS \\
NIRM : 950051053114120051

\vspace{1cm}
telah periksa, disetujui dan diuji oleh Pembimbing

\vspace{0.5cm}
Yogyakarta, \ldots\ Agustus 1999

\vspace{2.5cm}

\begin{minipage}{0.45\textwidth}
\centering
Pembimbing I

\vspace{2.5cm}
\rule{5cm}{0.4pt}\\
(Dr. Ir. FX. Nurwadji W., M.Sc.)
\end{minipage}%
\begin{minipage}{0.45\textwidth}
\centering
Pembimbing II

\vspace{2.5cm}
\rule{5cm}{0.4pt}\\
(Ir. Ch. Arief Sudibyo)
\end{minipage}

\vspace{2cm}
Disahkan oleh :\\
Ketua Program Studi Teknik Sipil

\vspace{2.5cm}
\rule{6cm}{0.4pt}\\
(Ir. Wiryawan Sardjono P., M.T.)

\end{center}
\clearpage

% Halaman Motto -- transcribed from Cover/motto.doc and Cover/edge.doc
\thispagestyle{empty}
\vspace*{6cm}
\begin{center}
\textit{Sepanjang apapun sebuah perjalanan, selalu dimulai dengan langkah pertama}
\end{center}
\clearpage

% "edge.doc" contains the dedication page laid out as a small repeated
% two-column card (title/author/NIRM/year) -- reproduced here as a simple
% dedication-style page rather than the original card layout.
\thispagestyle{empty}
\vspace*{4cm}
\begin{center}
{\large\bfseries TUGAS AKHIR}\\[0.3cm]
OPTIMASI BETON BERTULANG PADA STRUKTUR PORTAL RUANG

\vspace{1cm}
Yohan Naftali\\
No. Mahasiswa : 07712/TSS\\
NIRM : 950051053114120051

\vspace{1cm}
1999
\end{center}
\clearpage

% Kata Hantar (preface) -- transcribed verbatim from Cover/Kata hantar.doc
\chapter*{KATA HANTAR}
\addcontentsline{toc}{chapter}{KATA HANTAR}

Terima kasih kepada Allah Bapa di surga, karena dengan rahmat dan tuntunan
dari pelita Allah tugas akhir ini berhasil diselesaikan, segenap puji dan
syukur dipanjatkan demi kemuliaan nama Allah Yang Maha Pengasih.

Tugas akhir yang mengambil judul Optimasi Beton Bertulang Pada Struktur
Portal Ruang ini ditulis untuk memenuhi salah satu persyaratan memperoleh
gelar Sarjana Teknik dalam kurikulum Strata-1 Program Studi teknik Sipil,
Fakultas Teknik, Universitas Atma Jaya Yogyakarta.

Selama penulisan tugas akhir ini banyak doa, dukungan, dan bimbingan baik
secara moral maupun secara material yang diberikan oleh berbagai pihak
untuk memperlancar penulisan tugas akhir ini. Karena itu pada kesempatan
ini penulis ingin mengucapkan terima kasih yang sebesar-besarnya kepada:

\begin{enumerate}
\item Papa dan Mama tercinta atas doa dan restu.
\item Ir. John Tri Hatmoko, M.Sc., selaku Dekan Fakultas Teknik
\item Ir. Wiryawan Sardjono P., M.T., selaku Ketua Program Studi Teknik Sipil
\item Dr. Ir. FX. Nurwadji Wibowo, M.Sc., selaku Dosen Pembimbing I yang telah
      memberikan bimbingan dan mengarahkan penulis, sehingga penulis
      berkesempatan mengenal dunia optimasi dan pemrograman.
\item Ir. Ch. Arief Sudibyo, selaku Dosen Pembimbing II yang telah memberikan
      bimbingan kepada penulis
\item Segenap dosen dan karyawan Program Studi Teknik Sipil Universitas Atma
      Jaya.
\item Mutiara atas cinta dan kasihnya untuk penulis.
\item Bu Yenny, Pak Wagianto, dan teman-teman di Laboratorium Teknik
      Penyehatan.
\item Yerry, Bayu, Fika, Herman, Roy, Benny, Eko, Anwar dan semua teman yang
      tidak dapat disebutkan satu-persatu atas bantuan dan dukungannya.
\end{enumerate}

Semoga Allah Bapa memberikan imbalan atas budi baik yang telah ditanamkan
mereka.

\clearpage


\tableofcontents
\clearpage
\listoftables
\clearpage
\listoffigures
\clearpage

% Arti Lambang dan Singkatan -- transcribed verbatim from
% Daftar/Arti Lambang dan Singkatan.doc
\chapter*{ARTI LAMBANG DAN SINGKATAN}
\addcontentsline{toc}{chapter}{ARTI LAMBANG DAN SINGKATAN}

\begin{description}[style=nextline,leftmargin=2.2cm,labelwidth=2cm]
\item[$a$] Tinggi blok tegangan persegi ekivalen
\item[$AC$] Beban titik kumpul gabungan
\item[$AE$] Beban titik kumpul ekivalen
\item[$AFC$] Beban titik kumpul gabungan yang selaras dengan $DF$
\item[$AM_i$] Gaya di ujung batang $i$
\item[$AML$] Gaya ujung batang akibat beban
\item[$AMS_i$] Gaya jepit kedua ujung $i$ dalam arah sumbu struktur
\item[$AR$] Reaksi di titik kumpul yang dikekang
\item[$ARC$] Beban titik kumpul gabungan yang selaras dengan perpindahan pada titik kumpul yang dikekang
\item[$AS$] Luas tulangan pada daerah tarik
\item[$AS'$] Luas tulangan pada daerah tekan
\item[$b$] Lebar penampang beton
\item[$B$] Lebar kolom
\item[$c$] Jarak dari serat tekan terluar ke garis netral
\item[$CC$] Gaya internal pada daerah tekan beton
\item[$CS$] Gaya internal pada daerah tekan baja tulangan
\item[$d$] Jarak dari serat tekan terluar ke pusat tulangan tarik
\item[$d'$] Jarak dari serat tarik terluar ke pusat tulangan desak
\item[$dia$] Diameter tulangan pada kolom
\item[$diatldslp$] Diameter tulangan desak pada daerah lapangan balok
\item[$diatldstm$] Diameter tulangan desak pada daerah tumpuan balok
\item[$diatlgs$] Diameter tulangan geser
\item[$diatltrlp$] Diameter tulangan tarik pada daerah lapangan balok
\item[$diatltrtm$] Diameter tulangan tarik pada daerah tumpuan balok
\item[$DF$] Perpindahan titik kumpul bebas
\item[$DM_i$] Perpindahan di ujung batang $i$
\item[$DR$] Perpindahan titik kumpul yang dikekang
\item[$E$] Modulus elastisitas untuk gaya aksial
\item[$E_C$] Modulus elastisitas untuk beton
\item[$E_S$] Modulus elastisitas baja tulangan
\item[$f$] Fungsi sasaran
\item[$f_c'$] Kuat desak karakteristik beton
\item[$f_S'$] Tegangan yang terjadi pada baja tarik
\item[$f_y$] Kuat tarik baja tulangan
\item[$g_j$] Fungsi kendala pertaksamaan ke-$j$
\item[$G$] Modulus elastisitas untuk geser
\item[$h_j$] Fungsi kendala persamaan ke-$j$
\item[$H$] Tinggi kolom
\item[$I_x$] Momen inersia pada sumbu $x$
\item[$I_y$] Momen inersia pada sumbu $y$
\item[$I_z$] Momen inersia pada sumbu $z$
\item[$jaraktlgs$] Jarak tulangan geser
\item[$JVD$] Jumlah variabel desain
\item[$L$] Panjang bentang batang
\item[$m$] Jumlah batang
\item[$M$] Momen
\item[$M_n$] Momen tahanan nominal penampang
\item[$M_{ox}$] Momen desain efektif pada arah $x$
\item[$M_{oy}$] Momen desain efektif pada arah $y$
\item[$M_x$] Momen yang terjadi pada arah sumbu $x$
\item[$M_y$] Momen yang terjadi pada arah sumbu $y$
\item[$nB$] Jumlah balok pada struktur
\item[$nK$] Jumlah kolom pada struktur
\item[$ntldslp$] Jumlah tulangan desak pada daerah lapangan balok
\item[$ntldstm$] Jumlah tulangan desak pada daerah tumpuan balok
\item[$ntltrlp$] Jumlah tulangan tarik pada daerah lapangan balok
\item[$ntltrtm$] Jumlah tulangan tarik pada daerah tumpuan balok
\item[$N$] Jumlah tulangan pada salah satu sisi kolom
\item[$Ntitik$] Jumlah titik coba
\item[$P$] Gaya aksial
\item[$P_u$] Gaya aksial ultimit
\item[$SFF$] Gaya $AF$ akibat satu satuan perpindahan $DF$
\item[$SJ$] Matrik kekakuan titik terakit
\item[$SM$] Kekakuan batang dalam arah sumbu struktur
\item[$SMS_i$] Kekakuan batang untuk kedua ujung batang $i$ dalam arah sumbu struktur
\item[$SRF$] Reaksi $AR$ akibat satu perpindahan $DF$
\item[$T$] Gaya internal pada daerah tarik baja tulangan
\item[$x$] Variabel desain
\item[$X$] Ruang variabel desain
\item[$X_B$] Koordinat $x$ titik terbaik
\item[$X_G$] Koordinat $x$ titik \textit{good}
\item[$X_M$] Koordinat $x$ titik tengah
\item[$X_S$] Koordinat $x$ arah penelusuran
\item[$X_T$] Koordinat $x$ titik coba baru
\item[$X_W$] Koordinat $x$ titik terjelek
\item[$Y_B$] Koordinat $y$ titik terbaik
\item[$Y_G$] Koordinat $y$ titik \textit{good}
\item[$Y_M$] Koordinat $y$ titik tengah
\item[$Y_S$] Koordinat $y$ arah penelusuran
\item[$Y_T$] Koordinat $y$ titik coba baru
\item[$Y_W$] Koordinat $y$ titik terjelek
\item[$\alpha$] Koefisien biaksial
\item[$\beta_1$] Faktor reduksi tinggi blok tegangan
\item[$\varepsilon$] Faktor konvergensi
\item[$\varepsilon_s$] Regangan pada tulangan tarik
\item[$\varepsilon_s'$] Regangan pada tulangan desak
\item[$\phi$] Faktor reduksi kekuatan
\item[$\rho$] Rasio penulangan pada daerah tarik
\item[$\rho'$] Rasio penulangan pada daerah tekan
\item[$\rho_b$] Rasio imbang
\item[$\rho_{maks}$] Rasio penulangan maksimum yang diijinkan
\end{description}

\clearpage

% Intisari (abstract) -- transcribed verbatim from Isi/Intisari.doc
\chapter*{INTISARI}
\addcontentsline{toc}{chapter}{INTISARI}

\noindent\textbf{OPTIMASI BETON BERTULANG PADA STRUKTUR PORTAL RUANG},
Yohan Naftali, tahun 1999, Peminatan Program Studi Struktur, Program Studi
Teknik Sipil, Fakultas Teknik, Universitas Atma Jaya Yogyakarta.

\bigskip
\noindent
Desain struktur merupakan salah satu bagian dari keseluruhan proses
perencanaan bangunan. Semakin mahalnya harga material menyebabkan harga
struktur bangunan menjadi semakin mahal, oleh karena itu untuk menekan biaya
bangunan, perlu dicari dimensi batang sehingga memberikan harga struktur
yang paling murah (optimal), tanpa melanggar kendala-kendala yang ada. Salah
satu metoda untuk mendapatkan struktur yang optimal dapat dilakukan dengan
proses optimasi. Kasus optimasi yang terjadi dalam bidang teknik pada
umumnya berupa masalah yang tidak linier, untuk itu perlu digunakan metoda
optimasi yang dapat menyelesaikan masalah tidak linier, dari dua metoda
yaitu dengan bantuan kalkulus dan metoda numerik.

Beton bertulang merupakan material yang banyak digunakan untuk membuat
struktur bangunan karena material pembentuknya mudah didapat. Balok utama
pada struktur beton bertulang langsung ditumpu oleh kolom dan dianggap
menyatu secara kaku dengan kolom, sistem ini dikatakan sebagai sistem
portal. Portal ruang merupakan sistim portal yang dianalisa secara 3
dimensi. Portal ruang memiliki 6 perpindahan yang mungkin terjadi pada
setiap titik kumpulnya.

Harga struktur yang murah tetapi tidak melanggar kendala yang ada merupakan
fungsi sasaran yang dicari dalam tugas akhir ini. Variabel desainnya berupa
dimensi penampang beton dan tulangan baja yang digunakan. Balok dan kolom
berpenampang empat persegi panjang dan uniform sepanjang bentangnya. Hasil
tugas akhir ini berupa sebuah perangkat lunak untuk optimasi beton bertulang
pada struktur portal ruang. Kode program ditulis dengan bahasa C++ dan
dikompilasi menggunakan Borland C++.

Penyelesaian masalah optimasi beton bertulang pada struktur portal ruang
memakai metoda kalkulus terlalu rumit, oleh karena itu dalam tugas akhir ini
digunakan metoda optimasi polihedron fleksibel. Dalam menganalisa struktur
portal ruang digunakan metoda kekakuan. Balok dianalisa sebagai balok
berpenampang empat persegi panjang dan bertulangan rangkap, kolom dianalisa
sebagai kolom biaksial bertulangan simetri.

Hasil yang diperoleh dengan menggunakan program optimasi beton bertulang
pada struktur portal ruang yang dikembangkan dalam tugas akhir ini dapat
mengurangi harga struktur sekitar 9,9~\% dibandingkan hitungan struktur yang
dikerjakan tanpa dioptimasi. Dalam metoda polihedron fleksibel, semakin
banyak titik coba yang digunakan, harga struktur yang didapat menjadi
semakin murah, akan tetapi membutuhkan waktu proses optimasi yang semakin
lama. Pengembangan optimasi beton bertulang pada struktur portal ruang
merupakan hal yang masih luas dan menarik untuk dipelajari, terutama masalah
pendetailan antara sambungan balok kolom dan analisis struktur dinamis yang
belum dibahas dalam tugas akhir ini.

\clearpage


% =============================================================================
% BAGIAN ISI (MAIN MATTER) -- halaman diberi nomor angka Arab, mulai dari BAB I
% =============================================================================
\clearpage
\pagenumbering{arabic}
\setcounter{page}{1}

\chapter{PENDAHULUAN}
\label{bab:pendahuluan}

\section{Latar Belakang}

Desain struktur merupakan salah satu bagian dari keseluruhan proses
perencanaan bangunan. Proses desain dapat dibedakan menjadi dua tahap, tahap
pertama yaitu desain umum yang merupakan peninjauan secara garis besar
keputusan-keputusan desain, misalnya tata letak struktur, geometri atau
bentuk bangunan, dan material, sedangkan tahap kedua adalah desain secara
terinci yang antara lain meninjau tentang penentuan besar penampang balok
dan kolom, luas dan penempatan tulangan yang dibutuhkan pada struktur beton
bertulang dan sebagainya, untuk itu perlu dimodelkan struktur yang akan
didesain sehingga dapat diperoleh besaran gaya dan informasi perpindahan
yang diperlukan dalam menentukan luasan beton beserta penulangannya.

Untuk memodelkan suatu struktur bangunan yang baik dapat dimodelkan sebagai
portal ruang. Portal ruang merupakan suatu struktur yang pada tiap titiknya
memiliki enam perpindahan yang mungkin terjadi, sehingga bukan lagi dua atau
tiga seperti pada struktur rangka atau portal bidang.

Dalam menentukan dimensi balok dan kolom pada struktur beton bertulang harus
memperhatikan masalah kekuatan dan biaya, terutama pada saat ini dimana
harga material semakin mahal, banyak orang yang menunda bahkan membatalkan
proyeknya karena masalah biaya bahan bangunan yang mahal, oleh karena itu
dalam merencanakan struktur bangunan sangat perlu diperhatikan masalah
tersebut. Kekuatan yang dibutuhkan oleh suatu struktur beton bertulang dapat
dicapai dengan memberikan luasan penampang beton dan tulangan yang cukup.
Analisis perhitungan yang dilakukan diharapkan dapat memperoleh hasil yang
aman dan ekonomis. Untuk mendapatkan hasil yang paling murah (optimal) dapat
dicapai dengan menggunakan proses optimasi. Hasil yang didapat merupakan
hasil yang mempunyai harga struktur yang paling murah tetapi tetap masih
mampu mendukung beban struktur dengan aman.

Metoda optimasi semakin berkembang seiring dengan perkembangan komputer.
Dalam bidang teknik sipil, formulasi masalah optimasinya berupa formula
tidak linear, oleh karena itu metoda optimasi yang berkembang dalam bidang
teknik adalah metoda optimasi non-linear.

Dalam menyelesaikan masalah optimasi non-linear, dapat diselesaikan dengan
bantuan kalkulus yang memerlukan informasi turunan, hal ini dapat dilakukan
apabila masalah yang ada belum begitu komplek, untuk masalah yang sangat
komplek dikembangkan metoda yang idenya berasal dari alam, seperti metoda
algoritma genetik yang berasal dari teori Darwin ``Survival to the
fittest'', metoda \textit{simulated annealing} yang didasarkan pada proses
mengerasnya logam, metoda \textit{flexible polyhedron} yang dalam prosesnya
variabel desain akan membentuk suatu segi banyak yang makin lama semakin
mengecil. Dari banyak metoda di atas, sampai saat ini belum ada metoda yang
dapat digunakan untuk menyelesaikan segala kasus dengan memuaskan, karena
masing-masing metoda memiliki keunggulan dan kerugiannya sendiri-sendiri.

Oleh karena itu dalam tugas akhir ini, untuk melakukan optimasi beton
bertulang pada struktur portal ruang digunakan metoda polihedron fleksibel.

\section{Materi Tugas Akhir}

\subsection{Rumusan Masalah}

Proses optimasi adalah proses mencari nilai variabel desain yang memberikan
nilai optimal yang berupa nilai maksimum ataupun nilai minimum dari fungsi
sasaran tanpa melanggar kendala-kendala yang ada.

Secara umum masalah optimasi dapat dirumuskan sebagai:
\begin{equation}
\text{minimumkan } f(x),\quad x \in X
\label{eq:1-1}
\end{equation}
yang memenuhi kendala:
\begin{align}
g_j(x) &\leq 0, \quad j = 1,2,\ldots,n \label{eq:1-2}\\
h_j(x) &= 0, \quad j = 1,2,\ldots,m \label{eq:1-3}
\end{align}
pada persamaan di atas $x$ menyatakan variabel desain, $X$ menyatakan ruang
variabel desain, $f(x)$ adalah fungsi sasaran, $g(x)$ fungsi kendala
pertaksamaan, $h(x)$ fungsi kendala persamaan, $n$ jumlah kendala
pertaksamaan, $m$ jumlah kendala persamaan.

Harga struktur yang minimum merupakan fungsi sasaran dalam optimasi ini,
dengan variabel disainnya terdiri dari volume beton dan berat tulangan,
dalam hal ini tulangan terdiri dari tulangan tarik dan desak untuk balok,
tulangan geser untuk balok, tulangan kolom, tulangan geser untuk kolom.
Kendala pertaksamaan yang harus dipenuhi berupa momen yang terjadi pada
struktur harus lebih kecil atau sama dengan momen yang dapat didukung oleh
masing-masing batang balok maupun kolom, gaya geser yang terjadi pada
struktur harus lebih kecil atau sama dengan gaya geser yang didukung oleh
beton dan tulangan geser pada masing-masing batang balok maupun kolom, gaya
tekan maupun tarik yang terjadi harus lebih kecil atau sama dengan gaya
tekan atau tarik yang didukung oleh kolom, serta lendutan yang terjadi harus
lebih kecil atau sama dengan lendutan ijin.

Optimasi beton bertulang pada struktur portal ruang ini dapat dirumuskan
sebagai:

\medskip\noindent
\textbf{untuk balok}, minimumkan:
\begingroup\footnotesize
\begin{equation}
\begin{aligned}
f_1(&\mathit{lebar}, \mathit{tinggi}, \mathit{diatltrtm}, \mathit{ntltrtm}, \mathit{diatldstm}, \mathit{ntldstm},\\
    &\mathit{diatltrlp}, \mathit{ntltrlp}, \mathit{diatldslp}, \mathit{ntldslp}, \mathit{diatlgs}, \mathit{jaraktlgs}) =\\[2pt]
  &(\mathit{lebar} \times \mathit{tinggi} \times \mathit{panjang\ balok} \times \mathrm{Rp}(\text{harga beton})/\mathrm{m}^3)\\
 +&\left(\tfrac14 \pi\, \mathit{diatltrtm}^2 \times \mathit{ntltrtm} \times 2 \times 0.25 \times \mathit{panjang\ balok} \times (\text{berat tulangan})_{\mathrm{kg/m}^3} \times \mathrm{Rp}(\text{harga tulangan})/\mathrm{kg}\right)\\
 +&\left(\tfrac14 \pi\, \mathit{diatldstm}^2 \times \mathit{ntldstm} \times 2 \times 0.25 \times \mathit{panjang\ balok} \times (\text{berat tulangan})_{\mathrm{kg/m}^3} \times \mathrm{Rp}(\text{harga tulangan})/\mathrm{kg}\right)\\
 +&\left(\tfrac14 \pi\, \mathit{diatltrlp}^2 \times \mathit{ntltrlp} \times 0.5 \times \mathit{panjang\ balok} \times (\text{berat tulangan})_{\mathrm{kg/m}^3} \times \mathrm{Rp}(\text{harga tulangan})/\mathrm{kg}\right)\\
 +&\left(\tfrac14 \pi\, \mathit{diatldslp}^2 \times \mathit{ntldslp} \times 0.5 \times \mathit{panjang\ balok} \times (\text{berat tulangan})_{\mathrm{kg/m}^3} \times \mathrm{Rp}(\text{harga tulangan})/\mathrm{kg}\right)\\
 +&\Big(\tfrac14 \pi\, \mathit{diatlgs}^2 \times \big(2(\mathit{lebar} - 2\,\mathit{selimut\ beton}) + 2(\mathit{tinggi} - 2\,\mathit{selimut\ beton})\big)\\
   &\times (\mathit{panjang\ balok}/\mathit{jaraktlgs}) \times (\text{berat tulangan})_{\mathrm{kg/m}^3} \times \mathrm{Rp}(\text{harga tulangan})/\mathrm{kg}\Big)
\end{aligned}
\label{eq:1-4}
\end{equation}
\endgroup
yang memenuhi kendala pertaksamaan:
\begin{align}
M_{\text{intern}} &\geq M_{\text{extern}} \label{eq:1-5}\\
G_{\text{intern}} &\geq G_{\text{extern}} \label{eq:1-6}\\
\mathit{Lendutan} &\leq \mathit{Lendutan}_{\text{ijin}} \label{eq:1-7}
\end{align}
pada persamaan di atas $f_1$ menyatakan harga balok pada struktur, $lebar$
dan $tinggi$ menyatakan ukuran penampang balok. Variabel desainnya terdiri
atas $diatltrtm$ menyatakan diameter tulangan tarik pada daerah tumpuan,
$ntltrtm$ menyatakan jumlah tulangan tarik pada daerah tumpuan, $diatldstm$
menyatakan diameter tulangan desak pada daerah tumpuan, $ntldstm$
menyatakan jumlah tulangan desak pada daerah tumpuan, $diatltrlp$
menyatakan diameter tulangan tarik pada daerah lapangan, $ntltrlp$
menyatakan jumlah tulangan tarik pada daerah lapangan, $diatldslp$
menyatakan diameter tulangan desak pada daerah lapangan, $ntldslp$
menyatakan jumlah tulangan desak pada daerah lapangan, $diatlgs$ menyatakan
diameter tulangan geser, $jaraktlgs$ menyatakan jarak antar tulangan geser.
$M$ adalah momen dan $G$ adalah gaya geser.

\medskip\noindent
\textbf{untuk kolom}, minimumkan:
\begingroup\footnotesize
\begin{equation}
\begin{aligned}
f_2(B,H,dia,N,\mathit{diatlgskl},\mathit{jaraktlgskl}) =\;
  &(B \times H \times \mathit{tinggi\ kolom} \times \mathrm{Rp}(\text{harga beton})/\mathrm{m}^3)\\
 +&\left(\tfrac14 \pi\, dia^2 \times (4N-4) \times \mathit{tinggi\ kolom} \times (\text{berat tulangan})_{\mathrm{kg/m}^3} \times \mathrm{Rp}(\text{harga tulangan})/\mathrm{kg}\right)\\
 +&\Big(\tfrac14 \pi\, \mathit{diatlgskl}^2 \times \big(2(B - \mathit{selimut\ beton}) + 2(H - \mathit{selimut\ beton})\big)\\
   &\times (\mathit{tinggi\ kolom}/\mathit{jaraktlgskl}) \times (\text{berat tulangan})_{\mathrm{kg/m}^3} \times \mathrm{Rp}(\text{harga tulangan})/\mathrm{kg}\Big)
\end{aligned}
\label{eq:1-8}
\end{equation}
\endgroup
yang memenuhi kendala persamaan:
\begin{align}
M_{\text{intern}} &\geq M_{\text{extern}} \label{eq:1-9}\\
G_{\text{intern}} &\geq G_{\text{extern}} \label{eq:1-10}\\
P_{\text{intern}} &\geq P_{\text{extern}} \label{eq:1-11}
\end{align}
pada persamaan di atas $f_2$ menyatakan harga kolom pada struktur, dengan
variabel desain $B$ dan $H$ menyatakan ukuran penampang beton kolom, $dia$
menyatakan diameter tulangan, $N$ merupakan jumlah tulangan yang digunakan,
$diatlgskl$ merupakan diameter tulangan geser kolom, $jaraktlgskl$
menyatakan jarak antara tulangan geser pada kolom. $M$ adalah momen, $G$
adalah gaya geser, $P$ adalah gaya aksial pada kolom.

Harga struktur portal seluruhnya dapat dihitung dengan
\begin{equation}
f = \sum_{b=1}^{nB} f_{1b} + \sum_{k=1}^{nK} f_{2k}
\label{eq:1-12}
\end{equation}
pada persamaan di atas $f$ merupakan harga portal total, $f_{1b}$
menyatakan harga balok $b$, $f_{2k}$ menyatakan harga kolom $k$, $nB$
menyatakan jumlah balok, $nK$ menyatakan jumlah kolom.

Untuk menyelesaikan masalah tersebut akan dibuat program komputer yang
ditulis dalam bahasa pemrograman yang sudah banyak digunakan yaitu bahasa
C++ dan menggunakan data diskrit serta menggunakan metoda optimasi
polihedron fleksibel yang dapat digunakan untuk menangani masalah
non-linear.

\subsection{Batasan Masalah}

Dalam melakukan optimasi beton bertulang pada struktur portal ruang ini
perhitungan beton bertulang didasarkan pada peraturan yang berlaku di
Indonesia yaitu Tata Cara Perhitungan Struktur Beton Untuk Bangunan Gedung
SK SNI T-15-1991-03 (PU, 1991), sedangkan untuk melakukan perhitungan gaya
pada struktur digunakan metoda kekakuan (Weaver, 1980). Mengingat luasnya
permasalahan dan keterbatasan waktu, maka dalam tugas akhir ini diperlukan
pembatasan masalah, yaitu:
\begin{enumerate}[label=\alph*.]
\item Analisa struktur portal ruang memakai metoda kekakuan
\item Rigid-end offset diabaikan
\item Beban yang bekerja pada struktur adalah beban gravitasi menerus
\item Balok dan kolom berpenampang empat persegi panjang dan uniform
      sepanjang bentangnya
\item Tidak ada gelincir antara beton dengan tulangan baja
\item Perencanaan kolom biaksial menggunakan metoda kontur beban
\item Gaya puntir diabaikan
\item Perencanaan beton bertulang menggunakan analisis ultimit
\end{enumerate}

\section{Maksud dan Tujuan Penulisan}

Penulisan tugas akhir ini dimaksudkan untuk memenuhi syarat yudisium dalam
mencapai tingkat kesarjanaan Strata 1 (S1) pada Program Studi Teknik Sipil,
Fakultas Teknik, Universitas Atma Jaya Yogyakarta.

Selain itu penulisan tugas akhir ini juga bertujuan untuk mengetahui lebih
dalam tentang metoda-metoda optimasi yang sekarang ini mulai dirasakan
manfaatnya dan mulai dikembangkan, terutama metoda yang dapat digunakan
untuk melakukan optimasi biaya portal beton bertulang, yang memiliki
permasalahan yang kompleks.

Hasil dari penulisan tugas akhir ini adalah berupa program komputer yang
ditulis dengan bahasa C++ yang dapat dikompilasi menggunakan Borland C++
5.0 atau versi yang lebih tinggi pada sistem operasi Windows 32 bit menjadi
sebuah perangkat lunak.

\section{Kegunaan Penulisan}

Perangkat lunak yang dihasilkan dalam penulisan ini berguna untuk
merencanakan ukuran penampang balok dan kolom beserta penulangannya dengan
memasukkan informasi tentang bentuk geometri portal ruang yang akan
direncanakan.

Selain itu program yang dibuat ini dapat dijadikan dasar pengembangan
perangkat lunak untuk menyelesaikan masalah optimasi sehingga dapat
tercipta sebuah program yang lebih baik lagi, sehingga akan didapatkan
hasil yang lebih optimum lagi.

\section{Metoda Penulisan}

Dalam melakukan penulisan ini banyak langkah-langkah yang diambil sebelum
mendapatkan sebuah perangkat lunak yang siap digunakan untuk melakukan
optimasi.

Tahap pertama, dilakukan studi pustaka mengenai optimasi beton bertulang
pada portal ruang, metoda yang digunakan untuk menganalisis struktur
digunakan metoda kekakuan, beton dianalisis dengan analisis ultimit,
sedangkan kolom biaksial dianalisis menggunakan metoda kontur beban.

Tahap kedua adalah membuat suatu listing program untuk melakukan analisa
struktur menggunakan metoda kekakuan, juga membuat listing program untuk
melakukan analisa balok menggunakan analisis ultimit, serta listing program
untuk melakukan analisis kolom biaksial menggunakan metoda kontur beban.

Tahap ketiga adalah menggabungkan listing program yang telah diuji
validasinya masing-masing sebagai subrutin yang dipanggil oleh sebuah
fungsi utama, dalam fungsi utama itu dilakukan proses optimasi.

Tahap keempat adalah melakukan uji validasi dari perangkat lunak yang
digunakan dengan hasil yang pernah dihitung tanpa menggunakan proses
optimasi atau dengan perangkat lunak lain yang sudah jadi.

Tahap terakhir adalah mencoba perangkat lunak untuk memecahkan kasus-kasus
yang ada untuk melakukan penyusunan laporan tugas akhir secara lengkap.

\chapter{TINJAUAN PUSTAKA}
\label{bab:tinjauan}

\begin{quote}
\small\itshape
Catatan transkripsi: sebagian besar persamaan bernomor pada bab ini
(\eqref{eq:2-1}--\eqref{eq:2-25}) aslinya disisipkan sebagai objek OLE
\textit{MathType/Equation Editor} di dalam \texttt{BAB II.doc}, bukan sebagai
teks biasa, sehingga tidak dapat diambil langsung sebagai teks oleh
\texttt{antiword}. Persamaan-persamaan tersebut disusun ulang di sini
berdasarkan (a) teks naratif di sekitarnya yang tetap utuh, (b) kode sumber
C++ yang menjadi hasil tugas akhir ini (\texttt{Struktur.hpp},
\texttt{Solver.hpp}, \texttt{Balok.hpp}, \texttt{Kolom.hpp} pada
\texttt{Optimasi Beton/Source/}), dan (c) rujukan standar yang memang
disebut dalam teks (Weaver \& Gere 1980; SK SNI T-15-1991-03; Hulse \&
Mosley 1986). Lihat \texttt{LEGACY\_TRANSCRIPTION\_NOTES.md} untuk detail
metodologinya. Gambar 2-1 dan 2-2 (penampang dan diagram regangan) adalah
citra asli yang tidak dapat dipulihkan.
\end{quote}

\section{Struktur Portal Ruang}

\subsection{Tinjauan Umum}

Langkah pertama dalam mendesain struktur beton bertulang adalah memperoleh
informasi tentang momen, gaya geser, dan gaya aksial yang akan ditahan. Hal
ini biasanya melibatkan perhitungan numerik yang rumit, oleh karena itu
langkah terbaik adalah memprosesnya dengan program analisa struktur yang
telah dibuat dalam komputer (Hulse dan Mosley, 1986, halaman 17).

Dengan perkembangan komputer yang cukup pesat dan harganya relatif murah,
sekarang telah beredar program-program struktur, baik hasil buatan sendiri
maupun perangkat lunak yang diperdagangkan. Dengan perangkat ini suatu
analisis struktur dapat diselesaikan jauh lebih cepat, dan respon sistem
struktur dapat diperoleh. Saat ini analisis tiga dimensi telah umum
digunakan (Wahyudi dan Rahim, 1997, halaman 17).

Dalam menganalisa struktur rangka (\textit{framed structure}) dengan
bantuan komputer dapat digunakan metoda matriks, baik dengan metoda gaya
(\textit{flexibility}) maupun metoda kekakuan (\textit{stiffness}). Metoda
kekakuan atau lebih dikenal sebagai metoda perpindahan lebih sesuai untuk
diprogram dalam komputer, karena setelah model analisis struktur
ditentukan, pertimbangan teknis yang lebih lanjut tidak lagi diperlukan, di
sini terdapat perbedaan dengan metoda gaya, walaupun kedua metoda ini sama
bentuk matematisnya, pada metoda gaya besaran yang tidak diketahui adalah
gaya kelebihan (\textit{redundant}) yang dipilih secara sembarang,
sebaliknya dalam metoda kekakuan, yang tidak diketahui adalah perpindahan
titik kumpul yang tertentu secara otomatis. Jadi, jumlah yang tak diketahui
dalam metoda kekakuan sama dengan derajat ketidaktentuan
(\textit{indeterminacy}) kinematis struktur (Weaver dan Gere, 1980,
halaman 1).

Struktur rangka dibagi atas enam kategori: balok, rangka batang bidang,
rangka batang ruang, portal bidang, balok silang (\textit{grid}) dan portal
ruang, setiap struktur rangka terdiri dari batang-batang yang panjangnya
jauh lebih besar dibandingkan ukuran penampang lintangnya. Titik kumpul
struktur rangka adalah titik pertemuan batang-batang, termasuk tumpuan dan
ujung bebas suatu batang. Tumpuan bisa merupakan jepitan, sendi, atau
tumpuan rol, dalam kasus tertentu sambungan antara batang atau batang
dengan tumpuan bisa bersifat elastis (atau semi-tegar). Beban pada struktur
rangka bisa berupa gaya terpusat, beban tersebar, atau kopel (Weaver dan
Gere, 1980, halaman 2).

\subsection{Metoda Kekakuan}

Metoda kekakuan pertama dikembangkan dari superposisi gaya untuk koordinat
perpindahan bebas, kemudian metoda ini diformalisasi dan diperluas dengan
memakai matriks kesepadanan (\textit{compatibility}) dan konsep kerja maya
(\textit{virtual}). Pada pendekatan kerja maya, matriks kekakuan titik
kumpul (\textit{joint stiffness}) yang lengkap dirakit dengan perkalian
tiga matriks, walaupun versi formal tersebut menarik dan beraturan, tetapi
membutuhkan matriks kesepadanan yang besar dan jarang (\textit{sparse}).
Selain itu, elemen matriks ini tidak mudah dinilai dengan tepat. Baik cara
pembentukan matriks di atas ataupun perakitan dengan proses perkalian tidak
sesuai untuk diprogram pada komputer. Metodologi yang baik adalah menarik
ide dari kedua pendekatan tersebut dan menambahkan beberapa teknik yang
berorientasi pada komputer, sehingga didapatkan suatu cara yang disebut
metoda kekakuan langsung (\textit{direct stiffness method}) yang lebih
mudah untuk diprogram pada komputer (Weaver dan Gere, 1980, halaman 148).

Wahyudi dan Rahim (1997, halaman 17) menyebutkan prosedur umum metoda
kekakuan ataupun metoda elemen hingga dapat dibagi menjadi lima bagian
yaitu:
\begin{enumerate}[label=\arabic*.]
\item Penyusunan vektor beban dan matriks kekakuan elemen.
\item Penyusunan vektor beban dan matriks kekakuan global.
\item Penyelesaian persamaan kekakuan yang menghasilkan perpindahan global.
\item Perhitungan deformasi elemen dengan menggunakan transformasi koordinat.
\item Penentuan tegangan atau gaya-gaya elemen.
\end{enumerate}

Kunci penyederhanaan proses perakitan matriks kekakuan titik $SJ$ ialah
pemakaian matriks kekakuan batang untuk aksi dan perpindahan di kedua ujung
setiap batang. Jika perpindahan batang didasarkan pada koordinat struktur
global, maka ia akan berimpit dengan perpindahan titik kumpul. Dalam hal
ini semua komplikasi geometris harus ditangani secara setempat, dan
transfer informasi batang ke array struktur bersifat langsung, yaitu
matriks kekakuan dan vektor beban ekivalen dapat dirakit dengan penjumlahan
langsung sebagai ganti perkalian matriks. Dengan demikian, perakitan
matriks kekakuan titik untuk $m$ batang dapat dituliskan sebagai:
\begin{equation}
SJ = \sum_{i=1}^{m} SMS_i
\label{eq:2-1}
\end{equation}
dalam persamaan ini, simbol $SMS_i$ menyatakan matriks kekakuan batang
ke-$i$ dengan gaya ujung dan perpindahan (untuk kedua ujung) yang dalam
arah sumbu struktur. Agar dapat dijumlahkan, semua matriks kekakuan batang
harus diperluas ke ukuran yang sama seperti $SJ$ dengan menambahkan baris
dan kolom nol. Akan tetapi, operasi ini dapat dihindari dalam pembuatan
program dengan meletakan elemen dari $SMS_i$ ke lokasi yang sesuai pada
$SJ$. Demikian juga halnya, vektor beban ekivalen $AE$ dapat dibentuk dari
kontribusi batang sebagai berikut:
\begin{equation}
AE = \sum_{i=1}^{m} AMS_i
\label{eq:2-2}
\end{equation}
dalam hal ini $AMS_i$ adalah vektor gaya (aksi) jepit ujung dalam arah
sumbu struktur di kedua ujung batang $i$. Seperti pada $SMS_i$, secara
teoritis $AMS_i$ perlu diperbesar dengan bilangan nol agar dapat dijumlahkan
secara matriks, tetapi langkah ini bisa dihindari dalam program. Beban
titik tumpul ekivalen dalam persamaan \eqref{eq:2-2} kemudian dijumlahkan
dengan beban sebenarnya yang diberikan di titik kumpul untuk membentuk
vektor beban total atau gabungan.

Untuk memisahkan suku yang berkaitan dengan perpindahan bebas struktur dari
perpindahan pengekang tumpuan (\textit{support restraint}), matriks
kekakuan dan beban harus ditata ulang (\textit{rearranged}) dan disekat
(\textit{partitioned}). Penataan ulang suku-suku ini dapat dilakukan
setelah perakitan selesai. Dalam program komputer, penataan ulang ini lebih
mudah dilakukan bila informasi ditransfer dari array batang yang kecil ke
array struktur yang besar.

Setelah matriks kekakuan dan vektor beban ditata ulang (\textit{rearranged})
dan disekat-sekat (\textit{partitioned}), maka penyelesaian dasar untuk
perpindahan pada titik kumpul yang bebas akibat beban adalah:
\begin{equation}
DF = SFF^{-1} \, AFC
\label{eq:2-3}
\end{equation}
dalam persamaan ini $AFC$ adalah vektor beban titik kumpul gabungan
(sebenarnya dan ekivalen) yang selaras (\textit{correspond}) dengan $DF$.
Walaupun secara simbolis penyelesaian dalam persamaan \eqref{eq:2-3} lebih
mudah dituliskan sebagai hasil inversi matriks kekakuan $SFF$, tetapi
perhitungan $SFF^{-1}$ dalam program komputer hakekatnya tidak efisien.
Sebagai gantinya, matriks koefisien $SFF$ difaktorisasi dan penyelesaian
untuk $DF$ diperoleh dari sapuan (\textit{sweep}) kemuka dan belakang --
dalam program optimasi beton bertulang pada struktur portal ruang ini,
faktorisasi tersebut dilakukan dengan pendekatan Choleski yang dimodifikasi
oleh fungsi \texttt{banfac()} dan penyelesaian sapuan oleh fungsi
\texttt{bansol()} pada \texttt{Solver.hpp}.

Bagian lain yang hendak dicari adalah reaksi tumpuan dan gaya jepit ujung.
Bila hanya pengaruh beban yang diperhitungkan, persamaan untuk reaksi
adalah:
\begin{equation}
AR = SRF \, DF + ARC
\label{eq:2-4}
\end{equation}
Selain itu gaya ujung batang dituliskan sebagai:
\begin{equation}
AM_i = SM_i \, DM_i + AML_i
\label{eq:2-5}
\end{equation}
Banyak variasi dalam mengorganisasi metoda kekakuan bisa dilakukan untuk
pemrograman. Tahapan analisa di atas merupakan pendekatan teratur yang
memiliki beberapa segi menguntungkan bila berhadapan dengan masalah yang
besar dan rumit (Weaver dan Gere, 1980, halaman 151).

Untuk menganalisa struktur portal ruang digunakan matrik kekakuan batang
yang terdapat pada tabel 2-1 dan dijelaskan pada bagian \ref{sec:kekakuan-batang}.

\subsection{Kekakuan Batang Portal Ruang}
\label{sec:kekakuan-batang}

Untuk membentuk matriks kekakuan titik dengan proses penjumlahan yang
ditunjukan oleh persamaan \eqref{eq:2-1}, matriks kekakuan batang yang
lengkap dalam sumbu arah struktur perlu diturunkan dahulu. Selain itu,
matriks kekakuan batang $SM_i$ untuk sumbu arah batang diperlukan dalam
menghitung gaya ujung batang dengan persamaan \eqref{eq:2-5}. Kekakuan
batang terhadap sumbu batang selalu dapat diperoleh dan ditransformasi ke
sumbu struktur.

Portal ruang merupakan kasus yang paling umum untuk struktur rangka,
batang dapat mengalami sembarang kedua belas perpindahan, dengan demikian,
matriks kekakuan batang ($SM_i$) untuk struktur portal ruang berordo
$12 \times 12$ yang dapat dilihat dalam tabel 2-1, dan setiap kolom pada
matriks menyatakan aksi akibat salah satu perpindahan satuan.

Matriks kekakuan batang portal ruang diperlihatkan pada tabel 2-1, matriks
ini disekat untuk memisahkan bagian yang berkaitan dengan kedua ujung
batang, secara umum bentuknya adalah:
\begin{equation}
SM_i = \begin{bmatrix} SM_{jj} & SM_{jk} \\ SM_{kj} & SM_{kk} \end{bmatrix}
\label{eq:2-6}
\end{equation}
pada persamaan di atas subskrip $j$ dan $k$ pada submatriks di atas
menunjukan ujung batang (Weaver dan Gere, 1980, halaman 154).

% Tabel 2-1: Matriks Kekakuan Batang Portal Ruang
%
% The original "Tabel 2-1" was an embedded raster/vector image
% (Isi/Tabel 2.doc contains only a single [pic] anchor, no extractable
% text), so it could not be recovered from the .doc binary at all. The
% matrix below is reconstructed from the standard 12x12 local-axis space
% frame member stiffness matrix as published in Weaver, W. Jr. and Gere,
% J.M., "Matrix Analysis of Framed Structures" (cited throughout this
% chapter as the method's source), consistent with the degree-of-freedom
% ordering (u,v,w,theta_x,theta_y,theta_z per end) used in Struktur.hpp /
% Variabel.hpp (SM, SMRT arrays). VERIFY AGAINST THE ORIGINAL DOCUMENT
% before relying on the exact sign convention for structural design.
\begin{table}[htbp]
\centering
\caption{Matriks Kekakuan Batang Portal Ruang (direkonstruksi dari Weaver \& Gere, 1980)}
\label{tab:2-1}
\begin{align*}
a&=\frac{EA}{L}, &
t&=\frac{GJ}{L}\\
b_z&=\frac{12EI_z}{L^3}, & c_z&=\frac{6EI_z}{L^2}, & d_z&=\frac{4EI_z}{L}, & e_z&=\frac{2EI_z}{L}\\
b_y&=\frac{12EI_y}{L^3}, & c_y&=\frac{6EI_y}{L^2}, & d_y&=\frac{4EI_y}{L}, & e_y&=\frac{2EI_y}{L}
\end{align*}
\resizebox{\textwidth}{!}{$
SM_i =
\begin{array}{c|cccccccccccc}
 & u_j & v_j & w_j & \theta_{xj} & \theta_{yj} & \theta_{zj} & u_k & v_k & w_k & \theta_{xk} & \theta_{yk} & \theta_{zk} \\ \hline
u_j          &  a &    0 &    0 &  0 &    0 &    0 & -a &    0 &    0 &  0 &    0 &    0 \\
v_j          &  0 &  b_z &    0 &  0 &    0 &  c_z &  0 & -b_z &    0 &  0 &    0 &  c_z \\
w_j          &  0 &    0 &  b_y &  0 & -c_y &    0 &  0 &    0 & -b_y &  0 & -c_y &    0 \\
\theta_{xj}  &  0 &    0 &    0 &  t &    0 &    0 &  0 &    0 &    0 & -t &    0 &    0 \\
\theta_{yj}  &  0 &    0 & -c_y &  0 &  d_y &    0 &  0 &    0 &  c_y &  0 &  e_y &    0 \\
\theta_{zj}  &  0 &  c_z &    0 &  0 &    0 &  d_z &  0 & -c_z &    0 &  0 &    0 &  e_z \\
u_k          & -a &    0 &    0 &  0 &    0 &    0 &  a &    0 &    0 &  0 &    0 &    0 \\
v_k          &  0 & -b_z &    0 &  0 &    0 & -c_z &  0 &  b_z &    0 &  0 &    0 & -c_z \\
w_k          &  0 &    0 & -b_y &  0 &  c_y &    0 &  0 &    0 &  b_y &  0 &  c_y &    0 \\
\theta_{xk}  &  0 &    0 &    0 & -t &    0 &    0 &  0 &    0 &    0 &  t &    0 &    0 \\
\theta_{yk}  &  0 &    0 & -c_y &  0 &  e_y &    0 &  0 &    0 &  c_y &  0 &  d_y &    0 \\
\theta_{zk}  &  0 &  c_z &    0 &  0 &    0 &  e_z &  0 & -c_z &    0 &  0 &    0 &  d_z
\end{array}
$}
\end{table}


\section{Struktur Beton Bertulang}

\subsection{Tinjauan Umum}

Beton bertulang telah digunakan sebagai material bangunan di setiap negara.
Pada banyak negara seperti Amerika dan Kanada, beton bertulang merupakan
material struktur yang dominan dalam konstruksi bangunan. Keuntungan dari
penggunaan beton bertulang adalah tulangan baja dan bahan pembentuk beton
tersedia di banyak tempat, hanya dibutuhkan kemampuan yang relatif rendah
dalam konstruksi beton, dan dari segi ekonomi beton bertulang lebih
menguntungkan dibandingkan bahan konstruksi lainnya (MacGregor, 1997,
halaman 1).

Pada struktur beton bertulang, balok utama yang langsung ditumpu oleh
kolom dianggap menyatu secara kaku dengan kolom. Sistem kolom dan balok
induk seperti ini dikatakan sebagai sistem portal. Sistem portal telah
lama digunakan sebagai sistem bangunan yang dapat menahan beban vertikal
gravitasi dan lateral akibat gempa. Sistem ini memanfaatkan kekakuan
balok-balok utama dan kolom. Dengan demikian integritas antara balok utama
dan kolom harus mendapat perhatian dan pendetailan tersendiri, karena di
sekitar daerah ini timbul gaya geser dan momen yang besar yang dapat
menyebabkan retak atau patahnya penampang (Wahyudi dan Rahim, 1997,
halaman 14).

\subsection{Balok}

Balok adalah anggota struktur yang paling utama mendukung beban luar serta
berat sendirinya oleh momen dan gaya geser (MacGregor, 1997, halaman 85).

Dua hal utama yang dialami oleh suatu balok adalah kondisi tekan dan
tarik, yang antara lain karena adanya pengaruh lentur ataupun gaya
lateral. Padahal dari data percobaan diketahui bahwa kuat tarik beton
sangatlah kecil, kira-kira 10\%, dibandingkan kekuatan tekannya. Bahkan
dalam problema lentur, kuat tarik ini sering tidak diperhitungkan, sehingga
timbul usaha untuk memasang baja tulangan pada bagian tarik guna mengatasi
kelemahan beton tersebut, menghasilkan beton bertulang (Wahyudi dan Rahim,
1997, halaman 39).

Apabila penampang tersebut dikehendaki untuk menopang beban yang lebih
besar dari kapasitasnya, sedangkan di lain pihak seringkali pertimbangan
teknis pelaksanaan dan arsitektural membatasi dimensi balok, maka
diperlukan usaha-usaha lain untuk memperbesar kuat momen penampang balok
yang sudah tertentu dimensinya tersebut. Apabila hal demikian yang
dihadapi, SK-SNI T-15-1991-03 pasal 3.3.3 ayat 4 memperbolehkan penambahan
tulangan baja tarik lebih dari batas nilai $\rho_{maks}$ bersamaan dengan
penambahan tulangan baja di daerah tekan penampang balok. Hasilnya adalah
balok dengan penulangan rangkap dengan tulangan baja tarik dipasang di
daerah tarik dan tulangan tekan di daerah tekan. Pada keadaan demikian
berarti tulangan baja tekan bermanfaat untuk memperbesar kekuatan balok
(Dipohusodo, 1994, halaman 85).

Dalam praktek, sistem tulangan tunggal hampir tidak pernah dimanfaatkan
untuk balok, karena pemasangan batang tulangan tambahan di daerah tekan,
misalnya di tepi atas penampang tengah lapangan, akan mempermudah
pengaitan sengkang (\textit{stirrup}) (Wahyudi dan Rahim, 1997, halaman
60).

\gambarhilang{Penampang balok bertulangan rangkap: tulangan tekan pada
jarak $d'$ dari serat tekan terluar, tulangan tarik pada jarak (tinggi
efektif) $d$.}{Penampang Tulangan Rangkap}

\gambarhilang{Diagram regangan penampang: regangan serat tekan beton
$\varepsilon_c = 0{,}003$; garis netral pada jarak $c$ dari serat tekan
terluar.}{Diagram Regangan}

Pada gambar 2-1 tampak suatu penampang dengan tulangan tekan yang
ditempatkan dalam jarak $d'$ dari serat tekan terluar dan tulangan tarik
dalam jarak tinggi efektif $d$. Diagram regangan yang terjadi dilukiskan
dalam gambar 2-2, dengan regangan serat tekan beton dianggap telah mencapai
regangan maksimum 0,003. Garis netral terletak pada jarak $c$ yang belum
diketahui nilainya. Melalui perbandingan segitiga dalam diagram ini, dapat
ditentukan besar regangan masing-masing tulangan, yaitu:
\begin{align}
\varepsilon_s' &= 0{,}003\,\frac{c-d'}{c} \label{eq:2-7}\\
\varepsilon_s &= 0{,}003\,\frac{d-c}{c} \label{eq:2-8}
\end{align}
pada persamaan di atas $a = \beta_1 c$, dengan $a$ tinggi blok tegangan
persegi ekivalen dan $\beta_1$ faktor reduksi tinggi blok tegangan (SK SNI
T-15-1991-03 pasal 3.3.2 butir 7.(3): $\beta_1 = 0{,}85$ untuk $f_c' \le 30$
MPa, berkurang $0{,}008$ untuk setiap kenaikan 1 MPa kuat beton, dengan
batas bawah $\beta_1 \ge 0{,}65$ -- lihat \texttt{Balok::analisa()} pada
\texttt{Balok.hpp}).

Pertama-tama, kedua tulangan $AS$ dan $AS'$ diasumsikan telah mencapai
tegangan leleh $f_y$, sehingga keseimbangan gaya horisontal
$0{,}85f_c'ab + AS'f_y = AS f_y$ memberikan tinggi blok tegangan awal:
\begin{equation}
a = \frac{(AS-AS')f_y}{0{,}85\,f_c'\,b}
\label{eq:2-9}
\end{equation}
dan substitusi $a=\beta_1 c$ ke dalam persamaan \eqref{eq:2-7} dan
\eqref{eq:2-8} memberikan regangan pada trial pertama ini:
\begin{align}
\varepsilon_s' &= 0{,}003\,\frac{a-\beta_1 d'}{a} \label{eq:2-10}\\
\varepsilon_s &= 0{,}003\,\frac{\beta_1 d - a}{a} \label{eq:2-10b}
\end{align}
Bila kedua kondisi persamaan regangan tersebut terpenuhi ($\varepsilon_s
\ge \varepsilon_y$ dan $\varepsilon_s' \ge \varepsilon_y$, dengan
$\varepsilon_y = f_y/E_s$), tegangan pada baja tulangan menjadi
$f_s = f_s' = f_y$ dengan $f_s$ adalah tegangan pada baja tulangan tarik,
$f_s'$ adalah tegangan pada baja tulangan tekan. Diagram tegangan
internalnya dapat dianggap menjadi tiga bagian yaitu:
\begin{align}
\text{pada daerah tekan beton:} \quad & CC = 0{,}85\,f_c'\,a\,b \label{eq:2-11}\\
\text{pada daerah tekan baja tulangan:} \quad & CS = AS'\,f_y \label{eq:2-12}\\
\text{pada daerah tarik baja tulangan:} \quad & T = AS\,f_y \label{eq:2-13}
\end{align}
dengan $AS'$ dan $AS$ berturut-turut menyatakan luas baja tulangan tekan
dan tarik. Berdasarkan keseimbangan horisontal gaya internal $CC+CS=T$,
dihasilkan persamaan:
\begin{equation}
0{,}85\,f_c'\,a\,b + AS'\,f_y = AS\,f_y
\label{eq:2-14}
\end{equation}
dengan tinggi blok tegangan adalah persamaan \eqref{eq:2-9}, dan momen
tahanan nominal penampang adalah:
\begin{equation}
M_n = 0{,}85\,f_c'\,a\,b\left(d-\frac{a}{2}\right) + AS'\,f_y\,(d-d')
\label{eq:2-16}
\end{equation}
Momen desain diperoleh dengan memberikan faktor reduksi $\phi$ tertentu,
yang disarankan dalam SK SNI T-15-1991-03, ke dalam persamaan
\eqref{eq:2-16} (Wahyudi dan Rahim, 1997, halaman 61--63). Standar SK SNI
T-15-1991-03 pasal 2.2.3 ayat 2 memberikan faktor reduksi kekuatan $\phi$
untuk berbagai mekanisme, antara lain:
\begin{enumerate}[label=\alph*.]
\item Lentur tanpa beban aksial \dotfill $\phi = 0{,}80$
\item Geser dan puntir \dotfill $\phi = 0{,}60$
\item Tarik aksial, tanpa dan dengan lentur (sengkang) \dotfill $\phi = 0{,}65$
\item Tekan aksial, tanpa dan dengan lentur (spiral) \dotfill $\phi = 0{,}7$
\item Tumpuan pada beton \dotfill $\phi = 0{,}70$
\end{enumerate}
Dengan demikian dapat dinyatakan bahwa kuat momen $M_R$ (kapasitas momen)
sama dengan kuat momen ideal $M_n$ dikalikan faktor $\phi$ (Dipohusodo,
1994, halaman 41) -- inilah yang dihitung sebagai variabel \texttt{FMU}
pada \texttt{Balok::analisa()}.

Bila persamaan regangan \eqref{eq:2-7} dan \eqref{eq:2-8} tidak dipenuhi,
mungkin tegangan tulangan tekan ataupun tulangan tarik tidak mencapai
tegangan leleh materialnya, sehingga dalam hal ini persamaan tersebut
tidak berlaku lagi. Untuk dapat menentukan momen ketahanan penampang,
regangan aktual harus dihitung terlebih dahulu, dengan tegangan baja
sebagai fungsi regangan dikalikan modulus elastisitas $E_s = 2 \times
10^5$ MPa:
\begin{align}
f_s' &= E_s\,\varepsilon_s' = 600\,\frac{a-\beta_1 d'}{a} \label{eq:2-17a}\\
f_s &= E_s\,\varepsilon_s = 600\,\frac{\beta_1 d - a}{a} \label{eq:2-18a}
\end{align}
(dengan $f_s', f_s$ dibatasi pada rentang $-f_y \le f_s',f_s \le f_y$).
Dengan mensubtitusi persamaan \eqref{eq:2-17a} dan \eqref{eq:2-18a} ke
dalam keseimbangan gaya horisontal $0{,}85f_c'ab + AS'f_s' = AS f_s$
(bentuk umum persamaan \eqref{eq:2-14} untuk tulangan yang belum tentu
leleh), diperoleh persamaan kuadrat dalam $a$:
\begin{equation}
\bigl(0{,}85\,f_c'\,b\bigr)\,a^2 \;+\; \bigl(600\,AS' - AS\,f_y\bigr)\,a
\;-\; 600\,\beta_1\,AS'\,d' \;=\; 0
\label{eq:2-19}
\end{equation}
diselesaikan dengan rumus abc, diambil akar yang bernilai positif (fungsi
\texttt{balok::analisa()} pada \texttt{Balok.hpp} menghitung persamaan ini
lewat variabel bantu \texttt{ASOL,BSOL,CSOL,DSOL}). Sehingga, momen
tahanan penampang menjadi:
\begin{equation}
M_n = 0{,}85\,f_c'\,a\,b\left(d-\frac{a}{2}\right) + AS'\,f_s'\,(d-d')
\label{eq:2-20}
\end{equation}
dengan $a$ hasil penyelesaian persamaan \eqref{eq:2-19} dan $f_s'$ dari
persamaan \eqref{eq:2-17a}.

Momen desain dapat diperoleh dengan memberikan faktor reduksi tertentu,
yang disarankan dalam SK SNI T-15-1991-03, ke dalam persamaan
\eqref{eq:2-20} (Wahyudi dan Rahim, 1997, halaman 61--63).

Seperti pada tulangan tunggal, keruntuhan tarik atau tekan dapat pula
terjadi pada penampang tulangan rangkap. Bila hal tersebut terjadi,
keruntuhan tarik dengan pelelehan tulangan lebih disukai daripada
keruntuhan tekan dengan kehancuran beton yang mendadak. Keadaan ini dapat
dikendalikan dengan memberikan batasan tertentu pada tulangan terpasang.
SK SNI T-15-1991-03 pasal 3.3.3 menetapkan batasan:
\begin{equation}
\rho - \rho' \le 0{,}75\,\rho_b
\label{eq:2-21}
\end{equation}
dengan $\rho_b$ adalah rasio imbang (\textit{balance ratio}) baja tulangan
yang bersesuaian dengan penampang tulangan tunggal (Wahyudi dan Rahim,
1997, halaman 64). Untuk menentukan rasio penulangan keadaan seimbang
($\rho_b$) digunakan persamaan:
\begin{equation}
\rho_b = 0{,}85\,\beta_1\,\frac{f_c'}{f_y}\,\frac{600}{600+f_y}
\label{eq:2-22}
\end{equation}
dengan $f_c'$ dan $f_y$ dalam MPa, $\beta_1$ adalah konstanta yang
merupakan fungsi dari kelas kuat beton (Dipohusodo, 1994, halaman 37).
Standar SK SNI T-15-1991-03 menetapkan nilai $\beta_1$ diambil $0{,}85$
untuk $f_c' \le 30$ MPa, berkurang $0{,}008$ untuk setiap kenaikan 1 MPa
kuat beton, dan nilai tersebut tidak boleh kurang dari $0{,}65$
(Dipohusodo, 1994, halaman 31). Program mengimplementasikan persamaan
\eqref{eq:2-21}--\eqref{eq:2-22} sebagai \texttt{RHB} dan
\texttt{kendala\_rho\_b} pada \texttt{Balok.hpp}.

\subsection{Kolom Biaksial}

Kolom adalah batang tekan vertikal dari rangka (\textit{frame}) struktur
yang memikul beban dari balok. Kolom meneruskan beban-beban dari elevasi
atas ke elevasi yang lebih bawah hingga akhirnya sampai ke tanah melalui
fondasi. Karena kolom merupakan komponen tekan, maka keruntuhan pada suatu
kolom merupakan lokasi kritis yang dapat menyebabkan keruntuhan
(\textit{collapse}) lantai yang bersangkutan, dan juga merupakan batas
runtuh total (\textit{ultimate total collapse}) seluruh strukturnya. Oleh
karena itu dalam merencanakan kolom perlu lebih teliti, yaitu dengan
memberikan kekuatan cadangan yang lebih tinggi daripada yang dilakukan
pada balok dan elemen struktur horisontal lainnya, terlebih lagi
keruntuhan tekan tidak memberikan peringatan awal yang cukup jelas (Nawy,
1990).

SK SNI T-15-1991-03 (PU 1991) memberikan definisi, kolom adalah komponen
struktur bangunan yang tugas utamanya menyangga beban tekan aksial dengan
bagian tinggi yang tidak ditopang paling tidak tiga kali dimensi lateral
terkecil. Sedangkan komponen struktur yang menahan beban aksial vertikal
dengan rasio bagian tinggi dengan dimensi terkecil kurang dari tiga
disebut pedestal (Dipohusodo, 1994, halaman 287).

Bila eksentrisitas beban mempunyai harga kecil sehingga gaya aksial tekan
menjadi penentu, dan juga bila dikehendaki suatu kolom beton dengan
penampang lintang yang lebih kecil, maka umumnya distribusi tulangan
lebih baik dibuat merata di sekeliling sisi penampang tersebut. Untuk
distribusi tulangan semacam ini, baja tulangan yang terletak di bagian
tengah penampang akan menerima tegangan yang lebih kecil dibandingkan
tulangan lainnya. Ketika kapasitas ultimit kolom tersebut telah tercapai,
tegangan pada baja tulangan belum tentu mencapai tegangan lelehnya,
sedangkan baja tulangan yang berada di tepi kemungkinan besar sudah leleh
(Wahyudi dan Rahim, 1997, halaman 215).

Menurut Winter dan Nilson pada tahun 1994, mekanisme gaya aksial yang
bekerja bersamaan dengan lentur pada kedua arah dari sumbu utama penampang
terjadi pada kolom-kolom sudut bangunan, pada balok-balok yang membentuk
rangka dengan kolom dan menyalurkan momen-momen ujungnya ke kolom dalam
dua bidang yang tegak lurus. Keadaan yang sama juga dapat terjadi pada
kolom-kolom sebelah dalam, khususnya pada tata letak kolom yang tidak
teratur.

Kolom-kolom seperti pada sudut bangunan pada struktur rangka memikul gaya
aksial dan sekaligus momen dari dua sumbu aksis yang disebut kolom
biaksial. Untuk menyelesaikan masalah kolom biaksial pada kolom persegi,
dapat digunakan prosedur yang umum dipakai yaitu eksentrisitas biaksial,
$e_x$ dan $e_y$, digantikan dengan eksentrisitas uniaksial ekivalen
$e_{0x}$ atau $e_{0y}$, dan kolom didesain sebagai kolom uniaksial
(MacGregor, 1997, halaman 466).

Apabila eksentrisitas pada arah $x$ ($e_x$) dibagi dengan sisi pada arah
$x$ lebih besar daripada eksentrisitas pada arah $y$ ($e_y$) dibagi sisi
pada arah $y$, maka momen desain efektif dapat dihitung dengan:
\begin{equation}
M_{ox} = M_x + M_y\,\frac{1-\beta}{\beta}
\label{eq:2-23}
\end{equation}
Dengan cara yang sama apabila eksentrisitas pada arah $y$ dibagi dengan
sisi pada arah $y$ lebih besar daripada eksentrisitas arah $x$ dibagi sisi
arah $x$, maka momen desain efektifnya adalah:
\begin{equation}
M_{oy} = M_y + M_x\,\frac{1-\beta}{\beta}
\label{eq:2-24}
\end{equation}
dengan $M_x$ adalah momen yang terjadi pada arah $x$ dan $M_y$ adalah
momen pada arah $y$, $b$ adalah sisi efektif pada arah $y$, $h$ adalah
sisi efektif pada arah $x$, sedangkan $\beta$ adalah koefisien biaksial
yang untuk kolom persegi dengan tulangan simetris dapat dihitung dengan
rumus empiris (Hulse dan Mosley, 1986, halaman 163):
\begin{equation}
\beta = 0{,}3 + \frac{0{,}7}{0{,}6}\left(0{,}6 - \frac{P_u}{h\,b\,f_c'}\right)
\label{eq:2-25}
\end{equation}
dengan $\beta$ tidak boleh lebih kecil daripada $0{,}3$, dan $P_u$ adalah
gaya aksial. Persamaan \eqref{eq:2-23}--\eqref{eq:2-25} diimplementasikan
sebagai \texttt{MOX}, \texttt{MOY}, dan \texttt{beta} pada
\texttt{kolom::analisa()} di \texttt{Kolom.hpp}.

\section{Metoda Optimasi Struktur}

\subsection{Tinjauan Umum}

Penggunaan metoda optimasi dalam perencanaan struktur sebenarnya bukanlah
merupakan hal yang baru dan sudah banyak dikembangkan karena manfaatnya
yang banyak dirasakan. Pada tahun 1890 Maxwell mengemukakan beberapa teori
tentang desain yang rasional pada suatu struktur yang kemudian dikembangkan
lebih lanjut oleh Michell pada tahun 1904 (Wu, 1986, halaman 1). Beberapa
penelitian tentang optimasi struktur yang ditujukan untuk penggunaan
praktis telah dilakukan sekitar tahun 1940 dan 1950. Pada tahun 1960
Schmit mendemonstrasikan penggunaan teknik pemrograman non-linier untuk
desain struktur dan menyebutnya dengan istilah ``sintesa struktur'' (Wu,
1986, halaman 1). Komputer digital yang kemudian dibuat dan mampu untuk
memecahkan masalah numeris dalam skala besar telah memberikan momentum
yang besar untuk penelitian. Pada awal tahun 1970 optimasi struktur telah
menjadi sesuatu yang penting dalam berbagai aspek perancangan suatu
struktur (Wu, 1986, halaman 1).

Ada dua pendekatan utama dalam optimasi struktur. Pendekatan yang satu
menggunakan pemrograman matematika dan pendekatan yang lain menggunakan
metoda kriteria optimal. Kedua pendekatan ini masing-masing mempunyai
kelebihan dan kekurangan.

Setiap struktur rangka memiliki empat hal pokok dengan empat hal ini
dapat merupakan suatu variabel yang dapat diubah-ubah untuk mengoptimasi
struktur tersebut (variabel desain). Empat hal ini adalah ukuran elemen,
geometri struktur (posisi titik-titik kumpul), gambaran struktur
(bagaimana titik-titik kumpul tersebut dihubungkan oleh elemen-elemen),
dan material bangunan. Material bangunan biasanya ditentukan terlebih
dahulu. Persyaratan-persyaratan yang harus dipenuhi oleh struktur disebut
kendala. Dalam sebagian besar kasus, kendala berhubungan dengan kekuatan
dan defleksi struktur.

Dalam banyak kasus optimasi struktur yang telah dipecahkan sejauh ini,
variabel desain diasumsikan dapat berubah secara kontinyu. Akan tetapi
dalam desain yang sesungguhnya kadang-kadang dijumpai variabel desain
diskrit yang terbatas. Balok baja dan balok beton bertulang hanya ada
dalam ukuran standar. Untuk mengatasi masalah diskrit ini, beberapa usaha
telah dilakukan untuk memecahkan kasus-kasus tertentu, antara lain
menggunakan pendekatan \textit{branch-and-bound} (Wu, 1986, halaman 4).

Keberhasilan optimasi struktur sejauh ini masih terbatas. Pada satu sisi,
optimasi telah berhasil digunakan dalam perancangan berbagai jenis
struktur seperti \textit{overhead cranes}, bangunan-bangunan standard,
menara transmisi, girder dengan bentang pendek dan medium, dan bermacam-
macam kendaraan termasuk pesawat terbang, mobil, dan kapal. Pada sisi
lain, teknik-teknik optimasi masih kurang dapat diterapkan jika struktur
sangat besar atau jika kendala-kendalanya terlalu kompleks, belum lagi
jika ada pertimbangan-pertimbangan lain seperti standar produksi,
pengaruh estetika struktur, dan praktek-praktek konvensional dalam
industri (Wu, 1986, halaman 5).

Kesulitan utama dari penggunaan praktis optimasi struktur adalah lamanya
proses komputasi sehingga biaya perhitungan menjadi relatif tinggi. Akan
tetapi belakangan ini biaya perhitungan telah turun secara drastis karena
komputer pribadi makin cepat dan makin murah harganya. Dengan demikian
optimasi struktur yang sangat tergantung pada komputer akan semakin luas
aplikasinya.

Formulasi masalah untuk optimasi struktur dengan variabel desain yang
dapat berubah secara kontinyu dan menggunakan analisa struktur elastik
adalah program non-linier. Pendekatan untuk pemecahannya menggunakan
program linier sekuensial dan program non-linier (Wu, 1986, halaman 6).

Pendekatan linier sekuensialnya adalah melinierkan fungsi kendala
non-linier dan fungsi sasaran kemudian menyelesaikan masalah menggunakan
program linier berulang-ulang. Romstad dan Wang pada tahun 1967 dan 1968
memberikan penjelasan fisik tentang formulasi ini yaitu: yang tidak
diketahui adalah penambahan variabel desain dari suatu titik dalam daerah
layak yang diperoleh dari iterasi sebelumnya, dan kendala dari pendekatan
linier ini adalah perpindahan ijin dari titik kumpul dan gaya-gaya yang
terjadi pada elemen tidak boleh melebihi yang diijinkan. Pendekatan
non-liniernya adalah menyelesaikan masalah optimasi dengan menggunakan
program non-linier secara langsung. Akan tetapi persamaan analisa elastik
menunjukkan masalah yang besar sebab tidak ada algoritma pada program
non-linier yang dapat menangani kendala persamaan secara efisien. Teknik
yang biasa digunakan untuk menangani masalah ini adalah menggunakan
matriks untuk analisa struktur dalam algoritmanya (Wu, 1986, halaman 7).

Formulasi masalah untuk optimasi struktur dengan variabel desain yang
dapat berubah secara kontinyu dan menggunakan analisa struktur plastis
(analisa batas) adalah program linier (Wu, 1986, halaman 8). Masalah ini
dapat dipecahkan dengan mudah bahkan untuk struktur yang besar. Akan
tetapi analisa batas belum dipakai secara luas dalam desain struktur. Di
samping itu banyak kendala lain seperti defleksi tidak dapat dipenuhi
jika perancangan didasarkan pada analisa batas.

Optimasi struktur dengan sebagian atau seluruh variabel desain merupakan
nilai diskrit yang terbatas dan menggunakan analisa struktur elastis
memberikan formulasi masalahnya berupa masalah non-linier diskrit.
Beberapa pendekatan untuk memecahkan masalah ini telah diusulkan termasuk
program integer sekuensial, algoritma penelusuran arah variabel diskrit,
dan lain-lain (Wu, 1986, halaman 8).

Program integer sekuensial adalah ekstrapolasi dari program linier
sekuensial ke program diskrit. Dalam setiap iterasi, program linier
integer digunakan untuk menjamin penelusuran hanya berkisar dari satu
titik diskrit ke yang lain. Toakley pada tahun 1968 menemukan bahwa
pendekatan ini tidak efisien dan hasilnya tidak bisa diramalkan.
Reinschmidt pada tahun 1971 juga mempunyai kesimpulan yang sama meskipun
dia telah menunjukkan beberapa contoh kasus numerik yang dipecahkan
dengan cara ini, termasuk desain elastis rangka batang yang terdiri atas
9 elemen. Saglam pada tahun 1976 menggunakan pendekatan yang mirip dan
telah memecahkan beberapa contoh rangka batang.

Algoritma penelusuran arah variabel diskrit dikemukakan oleh Lai dan
Achenbach (1973, halaman 119--131). Struktur portal, kantilever, dan
pelat lantai dua lapis dengan kendala dinamis telah dioptimasi dengan
metoda ini.

Optimasi struktur dengan sebagian atau seluruh variabel desain merupakan
harga diskrit yang terbatas dan menggunakan analisa struktur plastis,
formulasi masalahnya adalah program integer atau program integer
campuran. Toakley (1968, halaman 1219) merumuskan masalah desain plastis
dengan variabel diskrit sebagai program integer dan telah dipecahkan. Fu
dan You pada tahun 1976 menggunakan metoda complex, tetapi mereka
menemukan metoda ini konvergen prematur, kadang-kadang jauh dari nilai
optimum. Levey dan Fu (1979, halaman 363--368) menggunakan metoda
complex-simplex. Lev (1977, halaman 365--371) mengusulkan algoritma
\textit{branch-and-bound} untuk memecahkan masalah diskrit dengan kendala
linier. Algoritma ini mengatasi keterbatasan dari beberapa algoritma
program integer yang lain.

Selama perkembangan optimasi struktur, para ilmuwan secara bertahap
mengembangkan banyak cara yang lain untuk merumuskan dan memecahkan
berbagai masalah sehingga banyak algoritma-algoritma yang diusulkan.
Sementara situasi ini dikenal secara umum sebagai sesuatu yang tidak dapat
dihindari dan harus dialami selama tahap awal perkembangan bidang
teknologi, ada juga saran untuk optimasi struktur dengan pendekatan yang
disatukan dan sistematis.

Dalam pendekatan sistematis yang digambarkan oleh Morris dan kawan-kawan
pada tahun 1982, algoritma dibagi dalam tiga bagian, tiap bagian mempunyai
fungsi yang terdefinisi secara baik, dan program komputer untuk tiap
bagian dapat dikembangkan dan diuji secara terpisah. Pendekatan ini
memperbolehkan penggunaan program secara berdiri sendiri yang kemudian
dapat dikumpulkan dengan cepat untuk persoalan desain yang spesifik.
Pendekatan ini dapat juga digunakan sebagai bahan perbandingan untuk
perkembangan langkah-langkah optimasi yang baru. Di samping itu,
pendekatan sistematik ini dapat dengan mudah diterapkan pada masalah yang
lebih kompleks di bidang teknik yang lain (Wu, 1986, halaman 10).

Menurut Kirsch pada tahun 1981 biasanya dalam suatu perencanaan terdiri
atas empat langkah yaitu:
\begin{enumerate}[label=\arabic*.]
\item Perumusan syarat-syarat fungsional, yaitu mencari dan merumuskan
      syarat-syarat fungsional yang dalam beberapa kasus tidak terlihat
      secara nyata.
\item Perencanaan dasar, misalnya pemilihan topologi, tipe struktur dan
      material.
\item Proses optimasi, yaitu untuk memperoleh kemungkinan perencanaan
      terbaik dengan kriteria, pertimbangan dan batas-batas yang ada.
\item Pendetailan, setelah seluruh penyajian optimasi, hasil yang didapat
      harus diperiksa dan dimodifikasi jika perlu.
\end{enumerate}

Berdasarkan berbagai kemajuan ilmu dan teknologi, perancangan struktur
bangunan harus direncanakan secara optimal yaitu struktur yang paling
ekonomis serta memenuhi segala persyaratan yang diinginkan. Oleh karena
itu perlu dikembangkan suatu sistem yang mampu menangani berbagai masalah
optimasi. Secara umum masalah optimasi ada empat jenis, yaitu:
\begin{enumerate}[label=\arabic*.]
\item Optimasi bentuk.
\item Optimasi topologi.
\item Optimasi geometri.
\item Optimasi ukuran penampang.
\end{enumerate}

Dalam metoda optimasi terdapat tiga besaran utama, yaitu:
\begin{enumerate}[label=\arabic*.]
\item \textbf{Variabel desain.} Besaran yang tidak berubah nilainya
      disebut parameter tetap, sedangkan yang nilai berubah selama proses
      optimasi disebut variabel desain. Variabel desain merupakan
      variabel yang dicari dalam masalah optimasi. Contohnya adalah
      ukuran komponen struktur dan geometri struktur. Data variabel
      desain ada dua macam, yaitu data diskrit dan data kontinu. Dalam
      beberapa kasus, khususnya optimasi bentuk dan geometri, variabel
      desain lebih sesuai dinyatakan sebagai variabel desain kontinu
      dibandingkan variabel diskrit.
\item \textbf{Fungsi kendala.} Fungsi kendala merupakan suatu fungsi yang
      memberikan batasan daerah layak dan daerah tak layak. Dalam bidang
      teknik terdapat dua macam kendala yaitu:
      \begin{enumerate}[label=\alph*.]
      \item Kendala rencana, yaitu kendala yang menentukan variabel
            desain selain yang memberikan batasan berdasarkan sifat.
            Kendala ini biasanya dapat dilihat secara nyata, misalnya
            batasan karena masalah fungsional, fabrikasi atau keindahan.
            Contoh kendala rencana adalah ketebalan plat, kemiringan
            atap.
      \item Kendala sifat, yaitu kendala yang didapat dari persyaratan
            sifat. Biasanya kendala ini tidak dapat terlihat secara nyata
            karena berhubungan dengan analisis struktur. Contoh kendala
            sifat adalah batas tegangan maksimum, perpindahan
            (\textit{displacements}) yang diijinkan, kekuatan tekuk.
      \end{enumerate}
\item \textbf{Fungsi sasaran.} Fungsi sasaran adalah suatu fungsi yang
      mengandung kriteria dari struktur yang diinginkan, misalnya
      struktur dengan berat paling ringan, dengan harga termurah, paling
      aman atau paling efisien. Pemilihan fungsi sasaran merupakan hal
      yang terpenting dalam proses optimasi agar dapat mencapai sasaran
      yang sebenarnya sedekat mungkin. Dalam beberapa situasi fungsi
      sasaran dapat terlihat jelas. Misalnya jika ingin mencari harga
      yang termurah maka fungsi sasarannya dapat diasumsikan ke dalam
      berat strukturnya. Namun terkadang sulit juga untuk menentukan
      harga yang sebenarnya dari sebuah konstruksi, misalnya struktur
      dengan berat paling ringan belum tentu yang termurah, karena
      biasanya masalah harga minimum akan termasuk juga harga bahan,
      fabrikasi, transportasi dan lain-lain. Masalah dalam optimasi dapat
      dibedakan menjadi dua, yaitu:
      \begin{enumerate}[label=\alph*.]
      \item Masalah optimasi linier merupakan dasar dari metoda optimasi
            secara matematis. Dalam masalah ini fungsi kendala dan fungsi
            sasarannya semuanya dinyatakan dalam fungsi linier. Fungsi
            kendala dapat berupa persamaan maupun pertidaksamaan, dan
            fungsi sasarannya berupa meminimumkan dan memaksimumkan.
      \item Masalah optimasi tak linier, yaitu bila fungsi kendala dan
            fungsi sasarannya tak linier.
      \end{enumerate}
\end{enumerate}

Masalah optimasi dalam bidang teknik, pada umumnya berupa masalah
optimasi tak linier. Masalah yang tak linier ini juga dapat dilinierkan,
tetapi memberikan hasil yang kurang akurat ditinjau dari segi teknik.
Oleh karena itu terpaksa diselesaikan memakai program tak linier yang
lebih sulit dipelajari dibandingkan program linier, karena memerlukan
matematika yang kompleks.

Penyelesaian masalah optimasi dapat dipakai dua cara yaitu:
\begin{enumerate}[label=\arabic*.]
\item \textbf{Metoda analisa.} Metoda ini menggunakan dasar teori
      matematika yang dibuat oleh Maxwell pada tahun 1890 dan Michell
      tahun pada 1904 dan memberikan hasil eksak namun hanya dapat
      digunakan untuk masalah optimasi yang sederhana saja karena pada
      beberapa masalah yang lebih kompleks pengolahan matematikanya
      sangat tidak sederhana (Wibowo, 1996, halaman 2).
\item \textbf{Metoda numerik.} Metoda optimasi numerik berkembang sejak
      ditemukannya komputer sebagai alat bantu hitung. \textit{Dynamic
      programming, integer programming, steepest descent, sequential
      unconstraint minimization technique, gradient projection}, dan
      \textit{penalty function} merupakan metoda optimasi numerik yang
      sering dipakai untuk menyelesaikan masalah optimasi di bidang
      sipil (Wibowo, 1996, halaman 3). Dalam metoda ini nilai yang akan
      dicari didekati dengan cara iterasi dan proses iterasi dihentikan
      apabila nilai yang dicari sudah cukup dekat dengan titik optimal
      yang sesungguhnya (Kirsch, 1981, halaman 5).
\end{enumerate}

Metoda-metoda tersebut di atas mempunyai kelemahan, yaitu mempunyai
peluang relatif besar untuk konvergen ke titik optimum lokal, bila
dipakai untuk menyelesaikan masalah optimasi dengan jumlah titik optimum
lebih dari satu. Hal ini dapat dimaklumi, karena penurunan perumusannya
berdasarkan asumsi fungsi konveks. Selain itu juga memerlukan analisis
struktur yang banyak, keharusan mengikutsertakan semua kendala dalam
model matematikanya dan modifikasi variabel yang kurang efisien (Wibowo,
1996, halaman 3).

\subsection{Metoda Optimasi Polihedron Fleksibel}
\label{sec:polihedron}

Metoda polihedron fleksibel merupakan metoda pengembangan dari metoda
\textit{simplex}, yang dikembangkan oleh Nelder-Mead (Haftka, 1991,
halaman 64), dasar pemikiran metoda \textit{simplex} adalah menurunkan
nilai fungsi sasaran secara kontinu dimulai dari suatu nilai fungsi awal
sampai mencapai nilai fungsi minimum terpenuhi (Haftka, 1991, halaman
64).

Metoda ini bermanfaat untuk mencari harga-harga ekstrem suatu fungsi
dengan banyak variabel, dengan turunan dari fungsi tersebut sulit untuk
dicari. Metoda polihedron fleksibel menggunakan pencerminan
(\textit{reflection}), ekspansi (\textit{expansion}) dan penyusutan
(\textit{contraction}) dalam melakukan penelusuran.

Polihedron Fleksibel menggunakan banyak titik coba, dengan jumlah titik
coba $N_{titik}$ minimum sama dengan jumlah variabel desain $JVD$
ditambah satu. $N_{titik} = JVD+1$ dipakai oleh Harb dan Mattews pada
tahun 1987. Box mengusulkan $N_{titik} = 2 \cdot JVD$, sedangkan
$N_{titik} = JVD+2$ diusulkan oleh Biles pada tahun 1983. Jumlah titik
coba makin banyak dimaksudkan untuk mengurangi terjadinya konvergen
prematur, tetapi memperlambat konvergensi.

Masalah meminimumkan fungsi sasaran $f=f(x,y)$, dengan dua variabel
desain $x$ dan $y$, dan memakai jumlah titik coba $N_{titik} = JVD+1 = 3$,
prosedurnya adalah sebagai berikut:
\begin{enumerate}[label=\arabic*.]
\item Tentukan atau acak tiga buah titik di dalam ruang variabel disain,
      kemudian hitung nilai fitnessnya. Titik terbaik disebut titik $B$
      (\textit{best}) dengan koordinat $(X_B,Y_B)$, titik terjelek
      disebut titik $W$ (\textit{worst}) dengan koordinat $(X_W,Y_W)$,
      sedang titik lainnya disebut titik $G$ (\textit{good}) dengan
      koordinat $(X_G,Y_G)$.
\item Hitung titik pusat $M$ dengan koordinat $(X_M,Y_M)$ yang didapat
      dengan merata-rata titik coba selain titik $W$, sehingga:
      \begin{align}
      X_M &= 0{,}5\,(X_B + X_G) \label{eq:2-26}\\
      Y_M &= 0{,}5\,(Y_B + Y_G) \label{eq:2-27}
      \end{align}
\item Tentukan arah penelusuran, yaitu garis yang menghubungkan titik $W$
      dan titik $M$ dan dinyatakan sebagai:
      \begin{align}
      X_S &= X_M - X_W \label{eq:2-28}\\
      Y_S &= Y_M - Y_W \label{eq:2-29}
      \end{align}
\item Cari titik coba baru dalam arah $S$ (\textit{search}) yang
      memberikan nilai fitness lebih baik daripada fitness di titik $W$,
      titik coba ini tidak boleh identik dengan titik $M$. Titik coba
      baru $T$ (\textit{try}) dengan koordinat $X_T,Y_T$ adalah:
      \begin{align}
      X_T &= X_W + \alpha\,X_S \label{eq:2-30}\\
      Y_T &= Y_W + \alpha\,Y_S \label{eq:2-31}
      \end{align}
      $\alpha$ adalah koefisien refleksi. Kalau tidak ditemukan titik
      coba baru yang lebih baik dari titik $W$ maka dilakukan penyusutan
      menuju $B$, besarnya penyusutan sama dengan setengah jaraknya
      terhadap titik $B$, sehingga titik $W$ baru adalah:
      \begin{align}
      X_W^{\text{baru}} &= 0{,}5\,(X_W + X_B) \label{eq:2-32}\\
      Y_W^{\text{baru}} &= 0{,}5\,(Y_W + Y_B) \label{eq:2-33}
      \end{align}
      dan titik $G$ baru adalah
      \begin{align}
      X_G^{\text{baru}} &= 0{,}5\,(X_W + X_G) \label{eq:2-34}\\
      Y_G^{\text{baru}} &= 0{,}5\,(Y_W + Y_G) \label{eq:2-35}
      \end{align}
      Kemudian diperiksa apakah sudah konvergen atau belum, kalau sudah
      maka berhenti, kalau belum ulangi ke langkah dua lagi.
      Pemeriksaan konvergensi dapat menggunakan persamaan-persamaan yang
      ada di bawah ini:
      \begin{align}
      |X_B - X_G| &\le \varepsilon \label{eq:2-36a}\\
      |X_B - X_W| &\le \varepsilon \label{eq:2-36b}\\
      |X_W - X_G| &\le \varepsilon \label{eq:2-36c}\\
      |Y_B - Y_G| &\le \varepsilon \label{eq:2-36d}\\
      |Y_B - Y_W| &\le \varepsilon \label{eq:2-36e}\\
      |Y_W - Y_G| &\le \varepsilon \label{eq:2-36f}
      \end{align}
      dengan $\varepsilon$ adalah suatu nilai konvergensi dan diambil
      sekecil mungkin, misalnya $0{,}0003$.
\end{enumerate}

Secara umum metoda ini membuat sebuah segi banyak dalam ruang variabelnya
yang terus diiterasi sehingga segi banyak itu makin lama makin mengecil,
dan akan didapatkan hasil yang optimum begitu segi banyak itu menjadi
sangat kecil sekali, yang ditentukan nilainya sebagai suatu nilai
konvergensi.

\chapter{OPTIMASI BETON BERTULANG PADA STRUKTUR PORTAL RUANG}
\label{bab:optimasi}

\begin{quote}
\small\itshape
Catatan transkripsi: persamaan kendala \eqref{eq:3-1}--\eqref{eq:3-14} pada
bab ini juga aslinya berupa objek \textit{Equation Native} (MathType).
Bentuk persis setiap persamaan di bawah ini disusun ulang dengan mencocokkan
teks naratif (yang menyebut nama variabel setiap kendala secara eksplisit)
terhadap implementasi C++ yang mengimplementasikannya persis
(\texttt{balok::analisa()}, \texttt{balok::sengkang\_balok()},
\texttt{lendutan()} pada \texttt{Balok.hpp}; \texttt{kolom::rho()},
\texttt{kolom::analisa()}, \texttt{kolom::jarak\_tulangan()},
\texttt{kolom::kelangsingan()}, \texttt{kolom::sengkang\_kolom()} pada
\texttt{Kolom.hpp}), sehingga isi matematisnya terjamin sama dengan yang
benar-benar dihitung oleh program, walaupun nomor urut persamaan yang
persis seperti naskah asli tidak dapat dipastikan 100\%.
\end{quote}

\section{Struktur Program}

Program optimasi beton bertulang pada struktur portal ruang ini ditulis
dengan menggunakan bahasa pemrograman C++. Kode Program ini ditulis pada
beberapa file yaitu file utama yang berekstensi \texttt{cpp}, dan file
pembantu berekstensi \texttt{hpp}. Di dalam file yang berekstensi
\texttt{cpp} terdapat suatu fungsi bernama \texttt{main()}, fungsi ini akan
dijalankan pertama kali oleh program, selanjutnya dari dalam fungsi ini
akan dipanggil fungsi-fungsi lainnya sehingga membentuk suatu program yang
utuh.

Struktur program optimasi beton bertulang pada struktur portal ruang ini
merupakan gabungan dari pemrograman terstruktur dan pemrograman
berorientasi obyek.

Pemrograman terstruktur memiliki kelebihan dalam pengorganisasian dalam
membuat kode-kode program sehingga mudah dibaca dan mudah untuk dipahami,
gagasan pemrograman terstruktur pertama kali dikemukakan oleh Profesor
Edsger W. Dijkstra dari University of Eindhoven Nederland pada tahun 1965
(Sutedjo dan Michael, 1997, halaman 1). Dalam tulisannya, beliau menyatakan
bahwa pernyataan GOTO seharusnya jangan dipergunakan, namun pernyataan ini
ditanggapi oleh HD. Mills yang beranggapan bahwa pemrograman terstruktur
tidak hanya dihubungkan dengan pernyataan GOTO saja, tetapi pada struktur
program itu sendiri yang menentukan apakah program tersebut terstruktur
atau tidak, baik menggunakan pernyataan GOTO maupun tidak.

Pemrograman berorientasi obyek diciptakan karena masih dirasakan adanya
keterbatasan pada bahasa pemrograman tradisional, konsep pemrograman
berorientasi obyek yaitu membagi masalah-masalah ke dalam obyek, sehingga
data dan fungsi-fungsi yang akan dioperasikan digabung menjadi satu
kesatuan disebut pengkapsulan (\textit{encapsulated}).

Dalam bidang ilmu teknik, pemrograman terstruktur lebih banyak digunakan
karena rumitnya masalah yang dikerjakan dan memerlukan kemudahan
pemeriksaan (\textit{debugging}), dalam program ini struktur program utama
menggunakan konsep pemrograman terstruktur, sedangkan konsep pemrograman
berorientasi obyek digunakan pada analisis balok dan analisis kolom. Kelas
balok dan kelas kolom akan dibangkitkan apabila diperlukan analisa balok
atau analisa kolom dan apabila tidak diperlukan lagi kelas-kelas itu
dihancurkan, sehingga penggunaan kelas ini dapat menghemat memori komputer.

File yang terdapat pada program optimasi beton bertulang pada struktur
portal ruang adalah:
\begin{description}[style=nextline]
\item[\texttt{ORCISF.CPP}] merupakan file yang memuat fungsi \texttt{main()}.
\item[\texttt{BORLANDC.HPP}] file yang memuat header pustaka yang dibuat oleh Borland.
\item[\texttt{HEADER.HPP}] berisi kumpulan header.
\item[\texttt{PROTO.HPP}] berisi prototipe program.
\item[\texttt{VARIABEL.HPP}] berisi deklarasi variabel global.
\item[\texttt{INOUT.HPP}] berisi kumpulan sub program penanganan masalah masukan dan keluaran.
\item[\texttt{STRUKTUR.HPP}] berisi kumpulan sub program analisa struktur dengan metoda kekakuan.
\item[\texttt{SOLVER.HPP}] berisi sub program pembantu penanganan matrik.
\item[\texttt{KOLOM.HPP}] berisi kelas kolom untuk analisa kolom biaksial bertulangan simetri.
\item[\texttt{BALOK.HPP}] berisi kelas balok untuk analisa balok persegi bertulangan rangkap.
\item[\texttt{ELEMEN.HPP}] berisi kumpulan sub program yang menangani elemen pada balok dan kolom.
\item[\texttt{POLYHEDRON.HPP}] berisi sub program untuk optimasi dengan menggunakan metoda polihedron fleksibel.
\item[\texttt{PENORMALAN.HPP}] berisi kumpulan sub program yang menangani perubahan variabel balok dan kolom menjadi variabel normal.
\item[\texttt{PENGACAKAN.HPP}] berisi sub program yang mengacak nilai variabel.
\item[\texttt{DISKRITISASI.HPP}] berisi kumpulan sub program yang merubah nilai kontinu menjadi diskrit.
\item[\texttt{KENDALA.HPP}] berisi sub program untuk menghitung besarnya kendala yang terjadi pada struktur.
\item[\texttt{TELUSUR.HPP}] berisi sub program untuk melakukan penelusuran variabel.
\item[\texttt{BARU.HPP}] berisi sub program untuk menentukan struktur baru dalam metoda optimasi.
\item[\texttt{PENGURUTAN.HPP}] berisi sub program untuk mengurutkan fitness struktur.
\item[\texttt{TAMPILAN.HPP}] berisi sub program yang menangani masalah tampilan.
\item[\texttt{CETAK.HPP}] berisi sub program yang menangani masalah pencetakan hasil akhir.
\end{description}

\section{Sub Program Masukan dan Keluaran}

\subsection{Sub Program Untuk Menampilkan Menu Utama}

Pada saat program dijalankan pertama kali akan muncul sebuah menu yang
berisi pilihan untuk menjalankan program, menu utama program optimasi
beton bertulang pada struktur portal ruang seperti diperlihatkan pada
gambar 3-1.

\begin{center}
\fbox{\parbox{0.85\textwidth}{\footnotesize\ttfamily\raggedright\noindent
*************************************************************************\\
\hphantom{x}      Program Optimasi Beton Bertulang Pada Struktur Portal Ruang\\
\hphantom{xxxxxxxxxxxxxxxxxxxxxxxxx}Oleh : Yohan Naftali\\
\hphantom{xxxxxxxxxxxxxxxxxxxxxxxxxxxxxx}7712/TS\\
*************************************************************************\\
\hphantom{xxxxxxxxxxxxxxxxxxxxxxxxxxxxxxxxxxxxxxxxxxxxxxxxxxxxx}Time : 23:09:48.10\\
\hphantom{x} 1. Input data awal ke file\\
\hphantom{x} 2. Input data beban ke file\\
\hphantom{x} 3. Melihat isi file input\\
\hphantom{x} 4. Mengoptimasi Struktur\\
\hphantom{x} 5. Keluar\\
\hphantom{x} Pilihan (1-5) = \_
}}
\captionof{figure}{Menu Utama Program (tangkapan layar konsol asli, direproduksi sebagai teks)}
\end{center}

Pilihan yang terdapat pada menu utama terdiri dari:
\begin{enumerate}[label=\arabic*.]
\item Input data awal ke file, untuk memasukkan data tentang informasi struktur.
\item Input data beban ke file, untuk memasukkan data beban yang bekerja pada struktur.
\item Melihat isi file input, untuk menampilkan data yang telah dimasukkan.
\item Mengoptimasi struktur, untuk memulai proses optimasi beton bertulang pada struktur portal ruang.
\item Keluar, untuk menghentikan program.
\end{enumerate}
Fungsi-fungsi yang digunakan untuk menampilkan menu utama yang terdapat
pada file \texttt{TAMPILAN.HPP} terdiri dari:
\begin{enumerate}[label=\arabic*.]
\item \texttt{void menu\_utama()}, fungsi ini dipanggil dari fungsi
      \texttt{void main()}, selanjutnya dari fungsi ini dipanggil
      fungsi-fungsi lainnya.
\item \texttt{void about()}, fungsi ini dipanggil untuk menampilkan
      informasi program pada layar.
\end{enumerate}

\subsection{Sub Program Memasukkan Data}

Untuk melakukan proses pemasukkan data (\textit{input}), dapat dilakukan
dengan dua cara yaitu:
\begin{enumerate}[label=\arabic*.]
\item Mengikuti panduan yang disediakan oleh program.
\item Mengedit teks ASCII dengan bantuan MS-DOS Editor atau Notepad.
\end{enumerate}
Memasukkan data melalui panduan program, maka pertama kali akan diminta
nama file generik (tanpa ekstensi), nama file ini digunakan program untuk
memberi nama file-file bentukan lainnya. Selanjutnya program meminta
data-data struktur, mengedit teks ASCII dengan bantuan ASCII editor lebih
mudah digunakan untuk mengedit data masukan yang telah ada.

Sebelum memulai pemasukan data ke komputer, maka data-data struktur harus
kita persiapkan terlebih dahulu supaya mempermudah proses pemasukan data,
data-data struktur yang perlu kita persiapkan adalah:
\begin{enumerate}[label=\arabic*.]
\item Nama struktur
\item Jumlah batang pada struktur
\item Jumlah titik kumpul
\item Jumlah pengekang tumpuan pada struktur
\item Jumlah titik kumpul yang dikekang
\item Koordinat titik kumpul
\item Alokasi elemen batang pada struktur
\item Pengekang titik kumpul
\item Kuat desak beton karakteristik
\item Kuat tarik tulangan utama
\item Kuat tarik tulangan sengkang
\item Beban yang mengenai struktur
\item Data lebar dan tinggi balok
\item Data sisi penampang kolom
\item Data diameter tulangan utama
\item Data diameter tulangan sengkang
\item Data jumlah tulangan
\item Data jarak antara tulangan sengkang
\end{enumerate}
Setelah data kita masukkan ke dalam komputer, selanjutnya diperiksa lagi
data yang telah kita masukkan untuk memastikan bahwa data yang telah
dimasukkan sudah benar, apabila ada kesalahan, maka data masukan dapat
dikoreksi dengan mengedit file masukan.

Struktur file masukan pada program optimasi beton bertulang pada struktur
rangka ini adalah:
\begin{enumerate}[label=\arabic*.]
\item Data umum struktur terdapat pada file \texttt{generik.inp}
\item Data beban yang bekerja pada struktur terdapat pada file \texttt{generik.bbn}
\item Data diskrit mengenai ukuran penampang batang terdapat pada file \texttt{generik.isd}
\item Data diskrit mengenai diameter tulangan utama terdapat pada file \texttt{generik.idl}
\item Data diskrit mengenai diameter tulangan sengkang terdapat pada file \texttt{generik.ids}
\item Data diskrit mengenai jumlah tulangan utama pada salah satu penampang terdapat pada file \texttt{generik.ijl}
\item Data diskrit mengenai jarak antara tulangan sengkang terdapat pada file \texttt{generik.ijs}
\end{enumerate}
Fungsi-fungsi yang digunakan dalam proses masukan adalah:
\begin{enumerate}[label=\arabic*.]
\item \texttt{void input\_data()}, terdapat pada file \texttt{INOUT.HPP}, untuk menampilkan menu masukan
\item \texttt{void input\_data\_umum()}, terdapat pada file \texttt{INOUT.HPP}, untuk memasukkan data umum struktur ke file
\item \texttt{void input\_data\_diskrit()}, terdapat pada file \texttt{INOUT.HPP}, untuk memasukkan data elemen batang
\item \texttt{void load\_data()}, terdapat pada file \texttt{PEMBEBANAN.HPP}, untuk memasukkan data beban yang bekerja pada struktur
\end{enumerate}

\subsection{Sub Program Untuk Membaca Data Masukan}

Setelah data masukan diproses oleh komputer, dan diletakkan ke dalam
file-file masukan, maka untuk dapat dilakukan proses optimasi maka data
masukan tersebut perlu dibaca ulang, fungsi-fungsi yang digunakan untuk
melakukan proses ini adalah:
\begin{enumerate}[label=\arabic*.]
\item \texttt{void baca\_data()}, terdapat pada file \texttt{INOUT.HPP}, untuk membaca data masukan
\item \texttt{void baca\_beban()}, terdapat pada file \texttt{PEMBEBANAN.HPP}, untuk membaca data pembebanan.
\end{enumerate}

\section{Sub Program Inti}

\subsection{Sub Program Analisa Struktur}

Untuk menganalisa struktur portal ruang digunakan metoda kekakuan yang
dikembangkan oleh Weaver dan Gere (1980). Secara umum proses dalam
analisa struktur adalah sebagai berikut:
\begin{enumerate}[label=\arabic*.]
\item Pembentukan matrik rotasi
\item Pembentukan matrik kekakuan batang
\item Pembentukan vektor beban
\item Perhitungan perpindahan pada titik kumpul
\item Perhitungan gaya dan reaksi pada struktur
\end{enumerate}
Sub program analisa struktur terdapat pada file \texttt{STRUKTUR.HPP},
\texttt{PEMBEBANAN.HPP}, dan \texttt{SOLVER.HPP}. File \texttt{STRUKTUR.HPP}
berisi fungsi-fungsi untuk membentuk matrik rotasi, pembentukan matrik
kekakuan batang, perhitungan perpindahan, dan perhitungan gaya dan reaksi
pada struktur. Fungsi untuk membentuk vektor beban terdapat pada file
\texttt{PEMBEBANAN.HPP}, sedangkan file \texttt{SOLVER.HPP} terdapat
fungsi \texttt{void banfac()}, yang digunakan untuk melakukan faktorisasi
matrik simetris berjalur dengan pendekatan Choleski yang dimodifikasi, dan
terdapat pula fungsi \texttt{void bansol()}, yang digunakan untuk
mengolah matrik berjalur.

\subsection{Sub Program Analisa Balok Persegi Panjang}

Sub program analisa balok persegi panjang terdapat pada file
\texttt{BALOK.HPP}. Dalam file \texttt{BALOK.HPP} terdapat sebuah kelas
yang bernama \texttt{class balok}. Kelas ini dibangkitkan untuk menghitung
kendala yang terdapat pada setiap balok pada struktur.

Persamaan perhitungan kendala-kendala balok yang terdapat pada kelas
balok adalah:

\begin{enumerate}[label=\arabic*.]
\item \textbf{Kendala rasio tulangan maksimum}
\begin{equation}
\mathit{kendala\_rho\_b} = \max\!\left(\frac{\rho}{\rho_b} - 1,\; 0\right)
\label{eq:3-1}
\end{equation}
pada persamaan di atas $\rho$ (\texttt{RHO}) adalah rasio penulangan,
$\mathit{kendala\_rho\_b}$ adalah kendala rasio tulangan maksimum yang
dibatasi oleh $\rho_b$ (\texttt{RHB}) yaitu rasio penulangan pada keadaan
luluh dikalikan $0{,}75$ (persamaan \eqref{eq:2-21}--\eqref{eq:2-22}).

\item \textbf{Kendala rasio penulangan minimum}
\begin{equation}
\mathit{kendala\_rho\_m} = \max\!\left(\frac{\rho_{min}}{\rho} - 1,\; 0\right),
\qquad \rho_{min} = \frac{1{,}4}{f_y}
\label{eq:3-2}
\end{equation}
pada persamaan di atas $\mathit{kendala\_rho\_m}$ adalah kendala rasio
penulangan minimum, $\rho_{min}$ (\texttt{RMIN}) adalah rasio penulangan
minimum.

\item \textbf{Kendala momen lentur}
\begin{equation}
\mathit{kendala\_M} = \max\!\left(\frac{M_u}{\phi M_n} - 1,\; 0\right)
\label{eq:3-3}
\end{equation}
pada persamaan di atas $\mathit{kendala\_M}$ merupakan kendala momen
lentur, $M_u$ adalah momen yang bekerja pada struktur, $\phi M_n$
(\texttt{FMU}) adalah momen yang dapat ditahan oleh balok (persamaan
\eqref{eq:2-16}/\eqref{eq:2-20}).

\item \textbf{Kendala lendutan}
\begin{equation}
\mathit{kendala\_lendutan} = \max\!\left(\frac{\mathit{Lendutan}}{\mathit{Lendutan}_{ijin}} - 1,\; 0\right)
\label{eq:3-4}
\end{equation}
pada persamaan di atas $\mathit{kendala\_lendutan}$ adalah kendala
lendutan, $\mathit{Lendutan}$ adalah lendutan yang terjadi,
$\mathit{Lendutan}_{ijin}$ adalah lendutan ijin maksimum (dihitung
menurut SK SNI T-15-1991-03 pasal 3.2.5, fungsi \texttt{lendutan()} pada
\texttt{Balok.hpp}).

\item \textbf{Kendala tulangan geser}
\begin{equation}
\mathit{kendala\_sb} = \max\!\left(\frac{Jarak\_S}{S_{maks}} - 1,\; 0\right)
\label{eq:3-5}
\end{equation}
pada persamaan di atas $\mathit{kendala\_sb}$ adalah kendala tulangan
geser, $Jarak\_S$ adalah jarak sengkang, $S_{maks}$ (\texttt{SmakS})
adalah jarak antar sengkang maksimum.
\end{enumerate}

\subsection{Sub Program Analisa Kolom Biaksial}

Sub program analisa kolom biaksial terdapat pada file \texttt{KOLOM.HPP}.
Dalam file \texttt{KOLOM.HPP} terdapat sebuah kelas yang bernama
\texttt{class kolom}. Kelas ini dibangkitkan untuk menghitung kendala
yang terdapat pada setiap kolom pada struktur.

Persamaan perhitungan kendala-kendala kolom yang terdapat pada kelas
kolom adalah:

\begin{enumerate}[label=\arabic*.]
\item \textbf{Kendala rasio tulangan maksimum}
\begin{equation}
\mathit{kendala\_r\_mak} = \max\!\left(\frac{\rho}{0{,}08} - 1,\; 0\right)
\label{eq:3-6}
\end{equation}
pada persamaan di atas $\rho$ (\texttt{RHO}) adalah rasio penulangan,
$\mathit{kendala\_r\_mak}$ adalah kendala rasio tulangan maksimum yang
dibatasi $0{,}08$ ($8\,\%$).

\item \textbf{Kendala rasio penulangan minimum}
\begin{equation}
\mathit{kendala\_r\_min} = \max\!\left(\frac{0{,}01}{\rho} - 1,\; 0\right)
\label{eq:3-7}
\end{equation}
pada persamaan di atas $\mathit{kendala\_r\_min}$ adalah kendala rasio
penulangan minimum yang dibatasi sebesar $0{,}01$ ($1\,\%$).

\item \textbf{Kendala gaya aksial minimum}
\begin{equation}
\mathit{kendala\_po} = \max\!\left(\frac{P_n}{P_o} - 1,\; 0\right),
\qquad P_o = \phi\left(0{,}85 f_c' \,\mathit{sisi}^2 + AS_{tot} f_y\right)
\label{eq:3-8}
\end{equation}
pada persamaan di atas $\mathit{kendala\_po}$ merupakan kendala gaya
aksial minimum atau disebut juga kendala eksentrisitas minimum, $P_n$
merupakan gaya aksial nominal dan $P_o$ merupakan batas gaya aksial yang
dapat dipikul oleh kolom dikalikan faktor reduksi kekuatan untuk kolom
berpengikat sengkang yaitu $0{,}8$.

\item \textbf{Kendala gaya aksial}
\begin{equation}
\mathit{kendala\_pn} = \max\!\left(\frac{P_n}{P_{n,coba}} - 1,\; 0\right)
\label{eq:3-9}
\end{equation}
pada persamaan di atas $\mathit{kendala\_pn}$ adalah kendala gaya aksial
nominal, $P_n$ merupakan gaya aksial nominal yang terjadi pada struktur,
$P_{n,coba}$ adalah gaya aksial yang mampu didukung oleh kolom,
$P_{n,coba}$ dicari bersama-sama $M_{n,coba}$ berdasarkan eksentrisitas
yang terjadi pada kolom dengan metoda \textit{false position} untuk
menentukan letak garis netral (fungsi \texttt{kolom::hitung\_kolom()} pada
\texttt{Kolom.hpp}).

\item \textbf{Kendala momen ekivalen biaksial}\\
Untuk $M_{nx} > M_{ny}$ digunakan persamaan:
\begin{equation}
\mathit{kendala\_mn} = \max\!\left(\frac{M_{ox}}{M_{n,coba}} - 1,\; 0\right)
\label{eq:3-11a}
\end{equation}
untuk $M_{nx} \le M_{ny}$ digunakan persamaan:
\begin{equation}
\mathit{kendala\_mn} = \max\!\left(\frac{M_{oy}}{M_{n,coba}} - 1,\; 0\right)
\label{eq:3-11b}
\end{equation}
pada persamaan di atas $\mathit{kendala\_mn}$ merupakan kendala momen
ekivalen biaksial, $M_{ox}$ merupakan momen ekivalen biaksial pada sumbu
$X$, $M_{oy}$ merupakan momen ekivalen biaksial pada sumbu $Y$ (persamaan
\eqref{eq:2-23}--\eqref{eq:2-24}), $M_{n,coba}$ merupakan momen nominal
pada arah sumbu biaksial yang ditinjau.

\item \textbf{Kendala jarak tulangan minimum}
\begin{equation}
\mathit{kendala\_tul} = \max\!\left(\frac{jarak_{min}}{\mathit{jarak\_antar\_tulangan}} - 1,\; 0\right)
\label{eq:3-12}
\end{equation}
pada persamaan di atas $\mathit{kendala\_tul}$ adalah kendala jarak antara
tulangan minimum, $jarak_{min}$ adalah jarak minimum yang diijinkan
($\max(1{,}5\,dia,\,40\text{ mm})$ dibalik: yang diambil adalah nilai
terkecil dari $1{,}5\,dia$ dan $40$ mm), $\mathit{jarak\_antar\_tulangan}$
adalah jarak antara tulangan utama.

\item \textbf{Kendala kelangsingan kolom}
\begin{equation}
\mathit{kendala\_kelangsingan} = \max\!\left(\frac{\mathit{rasio\_kelangsingan}}{22} - 1,\; 0\right),
\qquad \mathit{rasio\_kelangsingan} = \frac{K L}{\mathit{sisi}\sqrt{1/12}}
\label{eq:3-13}
\end{equation}
pada persamaan di atas $\mathit{kendala\_kelangsingan}$ merupakan kendala
kelangsingan pada kolom, $\mathit{rasio\_kelangsingan}$ merupakan rasio
kelangsingan pada kolom (dengan $K=0{,}5$), yang dibatasi minimal $22$
sesuai SK SNI T-15-1991-03 ($KL/r \le 22$).
\end{enumerate}

\subsection{Sub Program Metoda Optimasi Polihedron Fleksibel}

Untuk melakukan proses optimasi digunakan metoda polihedron fleksibel,
prosedur umum yang digunakan dalam program optimasi beton bertulang pada
struktur portal ruang dengan metoda polihedron fleksibel adalah:
\begin{enumerate}[label=\arabic*.]
\item Penentuan variabel untuk struktur generasi pertama, yang dilakukan
      dengan pengacakan, dalam program optimasi beton bertulang pada
      struktur portal ruang ini penentuan variabel terdapat pada file
      \texttt{PENGACAKAN.HPP} dalam melakukan pengacakan, sebuah struktur
      dipilih supaya menggunakan elemen yang paling besar nilainya
      supaya dalam salah satu struktur tersebut ada yang tidak memiliki
      kendala, hal ini dimaksudkan supaya dalam proses optimasi tidak
      terjebak kepada optimal lokal yang berkendala.
\item Pembangkitan generasi pertama, sesudah seluruh struktur
      dibangkitkan, nilai fitness, kendala, harga, dan variabel desain
      disimpan dalam array.
\item Seluruh struktur pada generasi pertama diurutkan berdasarkan
      fitnessnya, pengurutan ini terdapat pada file
      \texttt{PENGURUTAN.HPP}, metoda pengurutan yang digunakan adalah
      metoda \textit{Bubble Sort}, metoda ini berasal dari sifat
      gelembung air yang semakin ringan akan menuju ke permukaan air.
\item Penentuan arah baru untuk melakukan penelusuran, terdapat pada
      file \texttt{TELUSUR.HPP}.
\item Menggantikan struktur terjelek dengan struktur baru hasil
      penelusuran yang memiliki fitness lebih baik, apabila dalam
      penelusuran tidak ditemukan struktur yang lebih baik dari pada
      struktur terjelek, maka dilakukan penyusutan menuju struktur
      terbaik.
\item Selanjutnya dilakukan lagi pengurutan fitness, proses ini
      dilakukan berulang-ulang sampai pada suatu batas konvergensi.
\end{enumerate}

Formulasi untuk fungsi penalti yang digunakan dalam program optimasi
beton bertulang pada struktur portal ruang adalah \textit{exterior
penalty function}, proses penelusuran dimulai dari daerah tidak layak,
yang kemudian secara bertahap menuju daerah layak. Formulasi fungsi
penalti untuk menyelesaikan masalah meminimumkan $f(x)$ yang memenuhi
kendala $g_j(x) \le 0$, $j=1,2,3,\ldots,n_g$ adalah:
\begin{equation}
\phi(x) = f(x) + r \sum_{j=1}^{n_g} \max\bigl(g_j(x),\,0\bigr)
\label{eq:3-14}
\end{equation}
dalam persamaan di atas $\phi$ merupakan fungsi penalti, $f(x)$ merupakan
fungsi sasaran yang berupa harga struktur yang formulasinya terdapat pada
persamaan \eqref{eq:1-12}, $r$ parameter penalti, $n_g$ jumlah kendala,
dan $g_j(x)$ persamaan kendala yang telah dihitung dalam kelas balok dan
kelas kolom (jumlah seluruh $\mathit{kendala\_*}$ dari persamaan
\eqref{eq:3-1}--\eqref{eq:3-13}, dijumlahkan menjadi \texttt{kendala}
pada konstruktor \texttt{balok()}/\texttt{kolom()}). Nilai fitness yang
dipakai untuk mengurutkan struktur (\texttt{fitstr}, lihat
\texttt{Kendala.hpp}/\texttt{Polyhedron.hpp}) dihitung sebagai
$\mathit{finalti}/(\mathit{harga} + \mathit{finalti}\times\mathit{kendala})$,
bentuk kebalikan dari fungsi penalti \eqref{eq:3-14} sehingga struktur
termurah dan paling memenuhi kendala memberikan fitness tertinggi.

Nilai $r$ harus cukup besar agar mampu mengarahkan proses penelusuran
dekat dengan daerah layak, tetapi cukup kecil agar proses peminimuman
tidak mengalami kesulitan.

Dalam program optimasi beton bertulang pada struktur portal ruang
konvergensi ditentukan sebagai berikut:
\begin{enumerate}[label=\arabic*.]
\item Jumlah iterasi melebihi batas yang telah ditentukan
\item Terjadi penyusutan berturut yang menghasilkan fitness yang sama
      sebanyak jumlah struktur desain
\item Terdapat struktur berjumlah sama dengan jumlah variabel desain
      yang memiliki nilai fitness sama.
\end{enumerate}

\subsection{Sub Program Mencetak Hasil Akhir Optimasi}

Setelah suatu proses optimasi selesai, hasil optimasi berupa nilai-nilai
variabel desain terbaik, harga, kendala dan hasil perhitungan struktur
dicetak, sub program untuk mencetak hasil akhir tersebut ke dalam file
adalah \texttt{void cetak\_akhir()} yang terdapat pada file
\texttt{CETAK.HPP}.

File-file hasil keluaran yang dicetak oleh program optimasi beton
bertulang pada struktur portal ruang adalah:
\begin{enumerate}[label=\arabic*.]
\item \texttt{generik.opt} berisi dimensi balok dan kolom pada struktur
      yang paling optimal.
\item \texttt{generik.str} berisi hasil perhitungan gaya batang dengan
      metoda kekakuan.
\item \texttt{generik.kdl} berisi kendala yang terdapat pada struktur.
\item \texttt{generik.inf} berisi informasi data masukan.
\item \texttt{generik.his} berisi riwayat optimasi.
\end{enumerate}

\chapter{VALIDASI DAN APLIKASI PROGRAM OPTIMASI}
\label{bab:validasi}

\section{Validasi Program Optimasi}

Dalam melakukan validasi program optimasi, digunakan hasil perhitungan
beton bertulang pada struktur portal yang dikerjakan oleh Ir. Harsoyo
dalam bukunya berjudul Gedung Bertingkat II (Harsoyo, 1976), seri bina
bangunan konstruksi beton bertulang. Geometri bangunan yang dimodelkan
sebagai struktur portal bidang digambarkan dalam gambar 4-1.

\gambarhilang{Geometri struktur portal bidang menurut Harsoyo (1976) yang
dipakai sebagai kasus validasi.}{Geometri Struktur Portal Bidang}

Langkah pertama dalam analisa struktur setelah memodelkan struktur adalah
melakukan penomoran titik dan batang, dalam kasus ini terdapat 16 titik
kumpul dan 21 batang yang terdiri dari 12 kolom dan 9 balok. Penomoran
titik dan batang dapat dilihat dalam gambar 4-2.

\gambarhilang{Penomoran titik kumpul (1--16) dan batang (1--21, terdiri
dari 12 kolom dan 9 balok) pada struktur portal bidang validasi.}{Penomoran Titik Kumpul dan Batang Portal Bidang}

Beban yang bekerja pada struktur portal bidang tersebut berupa beban
merata pada balok dan beban akibat tekanan angin pada titik 5, 9, 13 ke
arah $-x$, besarnya beban merata pada batang diperlihatkan dalam tabel
4-1. Sedangkan besarnya beban akibat tekanan angin diperlihatkan dalam
tabel 4-2.

\begin{table}[htbp]
\centering
\caption{Beban Merata Pada Struktur Portal Bidang}
\label{tab:4-1}
\begin{tabular}{cccc}
\toprule
Nomor Batang & Beban Merata Total (t/m) & Berat Sendiri (t/m) & Beban Luar (t/m) \\
\midrule
13 & 4,100 & 0,360 & 3,740 \\
14 & 3,491 & 0,360 & 3,131 \\
15 & 4,100 & 0,360 & 3,740 \\
16 & 4,100 & 0,360 & 3,740 \\
17 & 3,491 & 0,360 & 3,131 \\
18 & 4,100 & 0,360 & 3,740 \\
19 & 1,762 & 0,300 & 1,462 \\
20 & 1,447 & 0,300 & 1,147 \\
21 & 1,762 & 0,300 & 1,462 \\
\bottomrule
\end{tabular}
\end{table}

\begin{table}[htbp]
\centering
\caption{Beban Tekanan Angin}
\label{tab:4-2}
\begin{tabular}{cc}
\toprule
Nomor Titik & Beban Tekanan Angin (kg) \\
\midrule
5  & 400 \\
9  & 800 \\
13 & 800 \\
\bottomrule
\end{tabular}
\end{table}

Dalam tugas akhir ini struktur portal bidang pada gambar 4-1 dioptimasi
dengan program optimasi beton bertulang pada struktur portal ruang, dan
hasil optimasinya dibandingkan dengan hasil yang telah dikerjakan oleh
Ir. Harsoyo D. (1976), semua data masukan dapat dilihat dalam lampiran
24, dengan menggunakan mutu beton 17,5 MPa dan mutu baja tulangan 240
MPa, tekuk pada kolom dianalisa. Perbandingan hasil akhir variabel desain
antara hasil yang dikerjakan oleh Ir. Harsoyo dan hasil optimasi program
optimasi beton bertulang pada struktur portal ruang disajikan pada tabel
4-3.

\begin{landscape}
\begingroup
\footnotesize
\begin{longtable}{clcccc}
\caption{Perbandingan Volume Beton dan Berat Besi}
\label{tab:4-3}\\
\toprule
\textbf{Batang} & \textbf{Variabel Desain} & \multicolumn{2}{c}{\textbf{Harsoyo (1976)}} & \multicolumn{2}{c}{\textbf{Naftali (1999)}} \\
\midrule
\endfirsthead
\multicolumn{6}{c}{\textit{(lanjutan Tabel 4-3)}}\\
\toprule
\textbf{Batang} & \textbf{Variabel Desain} & \multicolumn{2}{c}{\textbf{Harsoyo (1976)}} & \multicolumn{2}{c}{\textbf{Naftali (1999)}} \\
\midrule
\endhead
\bottomrule
\endfoot
\bottomrule
\endlastfoot
1 & Lebar / Tinggi & \multicolumn{2}{c}{0,30 m / 0,75 m} & \multicolumn{2}{c}{0,41 m / 0,41 m} \\
  & Volume beton & \multicolumn{2}{c}{0,9 m$^3$} & \multicolumn{2}{c}{0,6724 m$^3$} \\
  & Tulangan utama & \multicolumn{2}{c}{$10\phi25/4\phi16$ (187,929 kg)} & \multicolumn{2}{c}{$8\phi25$ (136,8158 kg)} \\
  & Tulangan sengkang & \multicolumn{2}{c}{$\phi8$--150} & \multicolumn{2}{c}{$\phi12$--300} \\
2 & Lebar / Tinggi & \multicolumn{2}{c}{0,30 m / 0,75 m} & \multicolumn{2}{c}{0,41 m / 0,41 m} \\
  & Volume beton & \multicolumn{2}{c}{0,9 m$^3$} & \multicolumn{2}{c}{0,6724 m$^3$} \\
  & Tulangan utama & \multicolumn{2}{c}{$10\phi25/4\phi16$ (187,929 kg)} & \multicolumn{2}{c}{$8\phi22$ (109,0117 kg)} \\
  & Tulangan sengkang & \multicolumn{2}{c}{$\phi8$--150} & \multicolumn{2}{c}{$\phi12$--300} \\
3 & Lebar / Tinggi & \multicolumn{2}{c}{0,30 m / 0,75 m} & \multicolumn{2}{c}{0,40 m / 0,40 m} \\
  & Volume beton & \multicolumn{2}{c}{0,9 m$^3$} & \multicolumn{2}{c}{0,64 m$^3$} \\
  & Tulangan utama & \multicolumn{2}{c}{$10\phi25/4\phi16$ (187,929 kg)} & \multicolumn{2}{c}{$8\phi22$ (101,2778 kg)} \\
  & Tulangan sengkang & \multicolumn{2}{c}{$\phi8$--150} & \multicolumn{2}{c}{$\phi8$--300} \\
4 & Lebar / Tinggi & \multicolumn{2}{c}{0,30 m / 0,75 m} & \multicolumn{2}{c}{0,43 m / 0,43 m} \\
  & Volume beton & \multicolumn{2}{c}{0,9 m$^3$} & \multicolumn{2}{c}{0,7396 m$^3$} \\
  & Tulangan utama & \multicolumn{2}{c}{$10\phi25/4\phi16$ (187,929 kg)} & \multicolumn{2}{c}{$8\phi19$ (85,6326 kg)} \\
  & Tulangan sengkang & \multicolumn{2}{c}{$\phi8$--150} & \multicolumn{2}{c}{$\phi12$--300} \\
5 & Lebar / Tinggi & \multicolumn{2}{c}{0,30 m / 0,60 m} & \multicolumn{2}{c}{0,44 m / 0,44 m} \\
  & Volume beton & \multicolumn{2}{c}{0,72 m$^3$} & \multicolumn{2}{c}{0,7744 m$^3$} \\
  & Tulangan utama & \multicolumn{2}{c}{$10\phi22/4\phi16$ (137,3734 kg)} & \multicolumn{2}{c}{$8\phi20$ (93,7609 kg)} \\
  & Tulangan sengkang & \multicolumn{2}{c}{$\phi8$--180} & \multicolumn{2}{c}{$\phi12$--300} \\
6 & Lebar / Tinggi & \multicolumn{2}{c}{0,30 m / 0,60 m} & \multicolumn{2}{c}{0,44 m / 0,44 m} \\
  & Volume beton & \multicolumn{2}{c}{0,72 m$^3$} & \multicolumn{2}{c}{0,7744 m$^3$} \\
  & Tulangan utama & \multicolumn{2}{c}{$8\phi22/4\phi16$ (113,0698 kg)} & \multicolumn{2}{c}{$8\phi19$ (77,8014 kg)} \\
  & Tulangan sengkang & \multicolumn{2}{c}{$\phi8$--180} & \multicolumn{2}{c}{$\phi8$--300} \\
7 & Lebar / Tinggi & \multicolumn{2}{c}{0,30 m / 0,60 m} & \multicolumn{2}{c}{0,42 m / 0,42 m} \\
  & Volume beton & \multicolumn{2}{c}{0,72 m$^3$} & \multicolumn{2}{c}{0,7056 m$^3$} \\
  & Tulangan utama & \multicolumn{2}{c}{$8\phi22/4\phi16$ (113,0698 kg)} & \multicolumn{2}{c}{$8\phi19$ (85,1948 kg)} \\
  & Tulangan sengkang & \multicolumn{2}{c}{$\phi8$--180} & \multicolumn{2}{c}{$\phi12$--300} \\
8 & Lebar / Tinggi & \multicolumn{2}{c}{0,30 m / 0,60 m} & \multicolumn{2}{c}{0,42 m / 0,42 m} \\
  & Volume beton & \multicolumn{2}{c}{0,72 m$^3$} & \multicolumn{2}{c}{0,7056 m$^3$} \\
  & Tulangan utama & \multicolumn{2}{c}{$10\phi22/4\phi16$ (137,3734 kg)} & \multicolumn{2}{c}{$8\phi22$ (109,4494 kg)} \\
  & Tulangan sengkang & \multicolumn{2}{c}{$\phi8$--180} & \multicolumn{2}{c}{$\phi12$--300} \\
9 & Lebar / Tinggi & \multicolumn{2}{c}{0,30 m / 0,60 m} & \multicolumn{2}{c}{0,41 m / 0,41 m} \\
  & Volume beton & \multicolumn{2}{c}{0,72 m$^3$} & \multicolumn{2}{c}{0,6724 m$^3$} \\
  & Tulangan utama & \multicolumn{2}{c}{$6\phi22/4\phi16$ (88,73483 kg)} & \multicolumn{2}{c}{$8\phi20$ (92,4476 kg)} \\
  & Tulangan sengkang & \multicolumn{2}{c}{$\phi8$--180} & \multicolumn{2}{c}{$\phi12$--300} \\
10 & Lebar / Tinggi & \multicolumn{2}{c}{0,30 m / 0,60 m} & \multicolumn{2}{c}{0,45 m / 0,45 m} \\
   & Volume beton & \multicolumn{2}{c}{0,72 m$^3$} & \multicolumn{2}{c}{0,81 m$^3$} \\
   & Tulangan utama & \multicolumn{2}{c}{$6\phi22/4\phi16$ (88,73483 kg)} & \multicolumn{2}{c}{$8\phi29$ (86,5081 kg)} \\
   & Tulangan sengkang & \multicolumn{2}{c}{$\phi8$--180} & \multicolumn{2}{c}{$\phi12$--300} \\
11 & Lebar / Tinggi & \multicolumn{2}{c}{0,30 m / 0,60 m} & \multicolumn{2}{c}{0,46 m / 0,46 m} \\
   & Volume beton & \multicolumn{2}{c}{0,72 m$^3$} & \multicolumn{2}{c}{0,8464 m$^3$} \\
   & Tulangan utama & \multicolumn{2}{c}{$6\phi22/4\phi16$ (88,73483 kg)} & \multicolumn{2}{c}{$8\phi28$ (161,6033 kg)} \\
   & Tulangan sengkang & \multicolumn{2}{c}{$\phi8$--180} & \multicolumn{2}{c}{$\phi8$--300} \\
12 & Lebar / Tinggi & \multicolumn{2}{c}{0,30 m / 0,60 m} & \multicolumn{2}{c}{0,44 m / 0,44 m} \\
   & Volume beton & \multicolumn{2}{c}{0,72 m$^3$} & \multicolumn{2}{c}{0,7744 m$^3$} \\
   & Tulangan utama & \multicolumn{2}{c}{$6\phi22/4\phi16$ (88,73483 kg)} & \multicolumn{2}{c}{$4\phi25$ (76,5066 kg)} \\
   & Tulangan sengkang & \multicolumn{2}{c}{$\phi8$--180} & \multicolumn{2}{c}{$\phi12$--300} \\
13 & Lebar / Tinggi & \multicolumn{2}{c}{0,25 m / 0,65 m} & \multicolumn{2}{c}{0,33 m / 0,57 m} \\
   & Volume beton & \multicolumn{2}{c}{0,975 m$^3$} & \multicolumn{2}{c}{1,1286 m$^3$} \\
   & Tulangan utama & \multicolumn{2}{c}{$3\phi25/3\phi25/2\phi25/4\phi25/3\phi25/3\phi25$ (228,126 kg)} & \multicolumn{2}{c}{$2\phi25/3\phi28/3\phi25/2\phi22/2\phi25/3\phi28$ (212,9059 kg)} \\
   & Tulangan sengkang & \multicolumn{2}{c}{$\phi8$--100} & \multicolumn{2}{c}{$\phi12$--130} \\
14 & Lebar / Tinggi & \multicolumn{2}{c}{0,25 m / 0,65 m} & \multicolumn{2}{c}{0,33 m / 0,52 m} \\
   & Volume beton & \multicolumn{2}{c}{0,65 m$^3$} & \multicolumn{2}{c}{0,6864 m$^3$} \\
   & Tulangan utama & \multicolumn{2}{c}{$3\phi25/3\phi25/3\phi25/2\phi25/3\phi25/3\phi25$ (143,9139 kg)} & \multicolumn{2}{c}{$2\phi22/3\phi22/3\phi25/3\phi20/2\phi22/3\phi22$ (111,5333 kg)} \\
   & Tulangan sengkang & \multicolumn{2}{c}{$\phi8$--100} & \multicolumn{2}{c}{$\phi12$--150} \\
15 & Lebar / Tinggi & \multicolumn{2}{c}{0,25 m / 0,65 m} & \multicolumn{2}{c}{0,34 m / 0,55 m} \\
   & Volume beton & \multicolumn{2}{c}{0,975 m$^3$} & \multicolumn{2}{c}{1,122 m$^3$} \\
   & Tulangan utama & \multicolumn{2}{c}{$3\phi25/3\phi25/2\phi25/4\phi25/3\phi25/3\phi25$ (228,126 kg)} & \multicolumn{2}{c}{$2\phi22/3\phi28/3\phi22/2\phi22/2\phi22/3\phi28$ (190,1989 kg)} \\
   & Tulangan sengkang & \multicolumn{2}{c}{$\phi8$--100} & \multicolumn{2}{c}{$\phi12$--140} \\
16 & Lebar / Tinggi & \multicolumn{2}{c}{0,25 m / 0,65 m} & \multicolumn{2}{c}{0,31 m / 0,55 m} \\
   & Volume beton & \multicolumn{2}{c}{0,975 m$^3$} & \multicolumn{2}{c}{1,023 m$^3$} \\
   & Tulangan utama & \multicolumn{2}{c}{$3\phi25/3\phi25/2\phi25/4\phi25/3\phi25/3\phi25$ (228,126 kg)} & \multicolumn{2}{c}{$2\phi28/3\phi28/3\phi25/2\phi28/2\phi28/3\phi28$ (225,0688 kg)} \\
   & Tulangan sengkang & \multicolumn{2}{c}{$\phi8$--100} & \multicolumn{2}{c}{$\phi12$--130} \\
17 & Lebar / Tinggi & \multicolumn{2}{c}{0,25 m / 0,65 m} & \multicolumn{2}{c}{0,33 m / 0,52 m} \\
   & Volume beton & \multicolumn{2}{c}{0,65 m$^3$} & \multicolumn{2}{c}{0,6864 m$^3$} \\
   & Tulangan utama & \multicolumn{2}{c}{$3\phi25/3\phi25/3\phi25/2\phi25/3\phi25/3\phi25$ (143,9139 kg)} & \multicolumn{2}{c}{$2\phi25/3\phi25/3\phi20/2\phi25/2\phi25/3\phi25$ (111,4915 kg)} \\
   & Tulangan sengkang & \multicolumn{2}{c}{$\phi8$--100} & \multicolumn{2}{c}{$\phi12$--170} \\
18 & Lebar / Tinggi & \multicolumn{2}{c}{0,25 m / 0,65 m} & \multicolumn{2}{c}{0,34 m / 0,54 m} \\
   & Volume beton & \multicolumn{2}{c}{0,975 m$^3$} & \multicolumn{2}{c}{1,1016 m$^3$} \\
   & Tulangan utama & \multicolumn{2}{c}{$3\phi25/3\phi25/2\phi25/4\phi25/3\phi25/3\phi25$ (228,126 kg)} & \multicolumn{2}{c}{$2\phi20/3\phi28/3\phi22/2\phi25/2\phi20/3\phi28$ (191,4206 kg)} \\
   & Tulangan sengkang & \multicolumn{2}{c}{$\phi8$--100} & \multicolumn{2}{c}{$\phi12$--140} \\
19 & Lebar / Tinggi & \multicolumn{2}{c}{0,25 m / 0,50 m} & \multicolumn{2}{c}{0,33 m / 0,49 m} \\
   & Volume beton & \multicolumn{2}{c}{0,75 m$^3$} & \multicolumn{2}{c}{0,9702 m$^3$} \\
   & Tulangan utama & \multicolumn{2}{c}{$3\phi19/4\phi19/2\phi19/4\phi19/3\phi19/4\phi19$ (179,4247 kg)} & \multicolumn{2}{c}{$2\phi25/3\phi28/3\phi22/2\phi19/2\phi25/3\phi28$ (172,2049 kg)} \\
   & Tulangan sengkang & \multicolumn{2}{c}{$\phi8$--100} & \multicolumn{2}{c}{$\phi12$--160} \\
20 & Lebar / Tinggi & \multicolumn{2}{c}{0,25 m / 0,50 m} & \multicolumn{2}{c}{0,33 m / 0,49 m} \\
   & Volume beton & \multicolumn{2}{c}{0,5 m$^3$} & \multicolumn{2}{c}{0,6468 m$^3$} \\
   & Tulangan utama & \multicolumn{2}{c}{$3\phi19/2\phi19/3\phi19/2\phi19/3\phi19/2\phi19$ (91,3503 kg)} & \multicolumn{2}{c}{$2\phi22/3\phi22/3\phi25/2\phi28/2\phi22/3\phi22$ (110,299 kg)} \\
   & Tulangan sengkang & \multicolumn{2}{c}{$\phi8$--100} & \multicolumn{2}{c}{$\phi12$--180} \\
21 & Lebar / Tinggi & \multicolumn{2}{c}{0,25 m / 0,50 m} & \multicolumn{2}{c}{0,33 m / 0,51 m} \\
   & Volume beton & \multicolumn{2}{c}{0,75 m$^3$} & \multicolumn{2}{c}{1,0098 m$^3$} \\
   & Tulangan utama & \multicolumn{2}{c}{$3\phi19/4\phi19/2\phi19/4\phi19/3\phi19/4\phi19$ (179,4247 kg)} & \multicolumn{2}{c}{$2\phi22/3\phi20/3\phi25/2\phi19/2\phi22/3\phi20$ (151,9952 kg)} \\
   & Tulangan sengkang & \multicolumn{2}{c}{$\phi8$--100} & \multicolumn{2}{c}{$\phi12$--170} \\
\midrule
\multicolumn{2}{l}{\textbf{Total Volume Beton}} & \multicolumn{2}{c}{\textbf{16,56 m$^3$}} & \multicolumn{2}{c}{\textbf{17,1624 m$^3$}} \\
\multicolumn{2}{l}{\textbf{Total Berat Besi}} & \multicolumn{2}{c}{\textbf{3113,5654 kg}} & \multicolumn{2}{c}{\textbf{2693,1281 kg}} \\
\end{longtable}
\endgroup
\end{landscape}

Selanjutnya dengan menetapkan harga satuan beton Rp.\ 250.000,00/m$^3$ dan
harga satuan besi Rp.\ 5.000,00/kg dapat dihitung harga struktur secara
keseluruhan seperti yang diperlihatkan dalam tabel 4-4, dalam hal ini
harga beton dan harga besi tersebut sudah termasuk biaya pemasangan
seperti upah pekerja, kawat pengikat dan papan acuan (\textit{form
work}).

Dari tabel 4-4 ternyata harga struktur secara keseluruhan dapat dikurangi
sebesar Rp.\ 1.951.587,00 atau sebesar 9,9~\% dengan melakukan optimasi.
Proses optimasi menggunakan 168 variabel desain dan 171 struktur desain
sebagai titik coba. Waktu proses optimasi yang diperlukan adalah 117
detik menggunakan komputer dengan prosesor Pentium Celeron 333 MHz dan
RAM 64 MB.

\begin{table}[htbp]
\centering
\caption{Perbandingan Harga Struktur}
\label{tab:4-4}
\begin{tabular}{lccccc}
\toprule
Material & Harga Satuan (Rp.) & \multicolumn{2}{c}{Harsoyo (1976)} & \multicolumn{2}{c}{Naftali (1999)} \\
\cmidrule(lr){3-4}\cmidrule(lr){5-6}
 & & Pemakaian & Harga (Rp.) & Pemakaian & Harga (Rp.) \\
\midrule
Beton & 250.000,00/m$^3$ & 16,56 m$^3$ & 4.140.000 & 17,1624 m$^3$ & 4.290.600 \\
Baja  & 5.000,00/kg      & 3113,5654 kg & 15.567.827 & 2693,1281 kg & 13.465.640 \\
\midrule
\textbf{Total} & & & \textbf{19.707.827} & & \textbf{17.756.240} \\
\bottomrule
\end{tabular}
\end{table}

\section{Aplikasi Program Optimasi}

Suatu struktur portal ruang seperti pada gambar 4-3 terbuat dari beton
bertulang memiliki data struktur sebagai berikut:
\begin{enumerate}[label=\arabic*.]
\item Jumlah batang 8 buah, terdiri dari 4 buah balok dan ditopang oleh 4
      buah kolom.
\item Tumpuan jepit pada titik kumpul 1, 2, 3, 4.
\item Kuat desak karakteristik beton = 30 MPa.
\item Kuat tarik baja tulangan memanjang = 400 MPa.
\item Kuat tarik baja tulangan sengkang = 240 MPa.
\item Selimut beton pada balok = 50 mm.
\item Selimut beton pada kolom = 50 mm.
\item Beban merata sebesar 35 kN/m pada semua balok.
\item Variabel desain yang digunakan terdapat pada lampiran.
\item Faktor penalti diambil $1\times10^{10}$.
\item Jumlah struktur desain sebagai titik coba diambil sama dengan
      jumlah variabel desain $+\,3$, jumlah variabel desain $\times\,2$
      dan jumlah variabel desain $\times\,3$.
\end{enumerate}

\gambarhilang{Geometri struktur portal ruang untuk kasus aplikasi
program.}{Geometri Struktur Portal Ruang}

Penomoran titik kumpul (\textit{joint}) dan penomoran elemen batang
diperlihatkan pada gambar 4-4, pada gambar 4-4 juga diperlihatkan
pembebanan pada struktur.

\gambarhilang{Penomoran titik kumpul dan elemen batang, serta pembebanan,
pada struktur portal ruang kasus aplikasi.}{Penomoran Titik dan Batang Portal Ruang}

Program dijalankan sebanyak 15 kali dengan 3 variasi jumlah struktur
desain, yaitu jumlah variabel desain $+\,3$, jumlah variabel desain
$\times\,2$, jumlah variabel desain $\times\,3$, masing-masing sebanyak
5 kali, dengan menggunakan komputer dengan prosesor Pentium Celeron
333 MHz dan RAM 64 MB, hasil optimasi disajikan dalam bentuk grafik pada
gambar 4-5.

\gambarhilang{Sumbu tegak: nilai fitness; sumbu datar: jumlah struktur
desain sebagai titik coba (tiga variasi: $JVD+3$, $JVD\times2$,
$JVD\times3$, masing-masing 5 pengulangan).}{Grafik Perbandingan Nilai Fitness}

Grafik pada gambar 4-5 menunjukkan bahwa semakin banyak struktur desain
yang digunakan, nilai fitness yang dicapai semakin tinggi, hal ini
berarti untuk mendapatkan hasil yang makin optimal dibutuhkan jumlah
struktur desain yang makin banyak, akan tetapi dengan bertambahnya jumlah
struktur desain, waktu yang dibutuhkan untuk proses optimasi juga
semakin banyak. Grafik perbandingan waktu proses optimasi yang
dibutuhkan oleh masing-masing penambahan jumlah variabel desain
ditunjukkan dalam grafik yang terdapat pada gambar 4-6.

\gambarhilang{Sumbu tegak: waktu proses optimasi (detik); sumbu datar:
jumlah struktur desain sebagai titik coba.}{Grafik Perbandingan Lama Proses Optimasi}

Hasil terbaik untuk kasus portal ruang pada gambar 4-3 dengan nilai
fitness $1174{,}59$ memerlukan waktu proses optimasi sebesar 74 detik.
Fungsi sasarannya yaitu harga struktur, dengan mengambil harga beton
Rp.\ 250.000,00/m$^3$ dan harga besi Rp.\ 5.000,00/kg adalah sebesar
Rp.\ 8.535.430,00, hasilnya ditunjukkan pada tabel 4-5, sedangkan
perhitungan rincian harga terdapat pada tabel 4-6.

\begingroup
\footnotesize
\begin{longtable}{clcc}
\caption{Dimensi Struktur Portal Ruang Hasil Optimasi}
\label{tab:4-5}\\
\toprule
\textbf{Batang} & \textbf{Variabel Desain} & \multicolumn{2}{c}{\textbf{Penggunaan Material}} \\
\midrule
\endfirsthead
\multicolumn{4}{c}{\textit{(lanjutan Tabel 4-5)}}\\
\toprule
\textbf{Batang} & \textbf{Variabel Desain} & \multicolumn{2}{c}{\textbf{Penggunaan Material}} \\
\midrule
\endhead
\bottomrule
\endfoot
\bottomrule
\endlastfoot
1 & Lebar / Tinggi & \multicolumn{2}{c}{0,45 m / 0,45 m (1,0125 m$^3$)} \\
  & Tulangan utama & \multicolumn{2}{c}{$8\phi19$ (102,4988 kg)} \\
  & Tulangan sengkang & \multicolumn{2}{c}{$\phi10$--300} \\
2 & Lebar / Tinggi & \multicolumn{2}{c}{0,50 m / 0,50 m (1,25 m$^3$)} \\
  & Tulangan utama & \multicolumn{2}{c}{$8\phi22$ (134,7477 kg)} \\
  & Tulangan sengkang & \multicolumn{2}{c}{$\phi10$--300} \\
3 & Lebar / Tinggi & \multicolumn{2}{c}{0,50 m / 0,50 m (1,25 m$^3$)} \\
  & Tulangan utama & \multicolumn{2}{c}{$8\phi28$ (208,6947 kg)} \\
  & Tulangan sengkang & \multicolumn{2}{c}{$\phi10$--300} \\
4 & Lebar / Tinggi & \multicolumn{2}{c}{0,45 m / 0,45 m (1,0125 m$^3$)} \\
  & Tulangan utama & \multicolumn{2}{c}{$8\phi22$ (127,9512 kg)} \\
  & Tulangan sengkang & \multicolumn{2}{c}{$\phi8$--300} \\
5 & Lebar / Tinggi & \multicolumn{2}{c}{0,25 m / 0,60 m (0,9 m$^3$)} \\
  & Tulangan utama & \multicolumn{2}{c}{$4\phi25/3\phi19/4\phi19/3\phi19/4\phi25/3\phi19$ (164,686 kg)} \\
  & Tulangan sengkang & \multicolumn{2}{c}{$\phi10$--160} \\
6 & Lebar / Tinggi & \multicolumn{2}{c}{0,30 m / 0,60 m (1,08 m$^3$)} \\
  & Tulangan utama & \multicolumn{2}{c}{$2\phi25/4\phi22/4\phi22/3\phi29/2\phi25/4\phi22$ (195,3126 kg)} \\
  & Tulangan sengkang & \multicolumn{2}{c}{$\phi10$--160} \\
7 & Lebar / Tinggi & \multicolumn{2}{c}{0,35 m / 0,50 m (1,05 m$^3$)} \\
  & Tulangan utama & \multicolumn{2}{c}{$2\phi19/3\phi25/4\phi28/2\phi22/2\phi19/3\phi25$ (171,1133 kg)} \\
  & Tulangan sengkang & \multicolumn{2}{c}{$\phi10$--160} \\
8 & Lebar / Tinggi & \multicolumn{2}{c}{0,30 m / 0,60 m (1,08 m$^3$)} \\
  & Tulangan utama & \multicolumn{2}{c}{$2\phi28/2\phi22/2\phi28/3\phi29/2\phi28/2\phi22$ (170,3307 kg)} \\
  & Tulangan sengkang & \multicolumn{2}{c}{$\phi10$--180} \\
\midrule
\multicolumn{2}{l}{\textbf{Total Volume Beton}} & \multicolumn{2}{c}{\textbf{8,635 m$^3$}} \\
\multicolumn{2}{l}{\textbf{Total Berat Besi}} & \multicolumn{2}{c}{\textbf{1275,336 kg}} \\
\end{longtable}
\endgroup

\begin{table}[htbp]
\centering
\caption{Rincian Harga Struktur Portal Ruang}
\label{tab:4-6}
\begin{tabular}{lccc}
\toprule
Material & Harga Satuan (Rp.) & Pemakaian Material & Harga (Rp.) \\
\midrule
Beton & 250.000,00/m$^3$ & 8,635 m$^3$   & 2.158.750 \\
Baja  & 5.000,00/kg      & 1275,336 kg   & 6.376.680 \\
\midrule
\textbf{Total} & & & \textbf{8.535.430} \\
\bottomrule
\end{tabular}
\end{table}

\chapter{KESIMPULAN DAN SARAN}
\label{bab:kesimpulan}

\section{Kesimpulan}

Dari tugas akhir ini dapat disimpulkan bahwa:
\begin{enumerate}[label=\arabic*.]
\item Program komputer optimasi beton bertulang pada struktur portal
      ruang yang dikembangkan dalam tugas akhir ini dengan menggunakan
      metoda optimasi polihedron fleksibel berhasil mengurangi harga
      struktur sekitar 9,9~\%.
\item Bertambahnya jumlah struktur desain sebagai titik coba
      menghasilkan struktur yang semakin murah, tetapi memerlukan waktu
      proses optimasi yang semakin lama.
\item Untuk jumlah variabel yang relatif besar, hasil yang didapatkan
      oleh metoda optimasi polihedron fleksibel kurang memuaskan, metoda
      polihedron fleksibel cukup baik hanya untuk kasus yang memiliki
      jumlah variabel desain yang relatif kecil.
\end{enumerate}

\section{Saran}

Untuk menyelesaikan masalah optimasi yang cukup rumit seperti optimasi
beton bertulang pada struktur portal ruang, diperlukan metoda optimasi
yang lebih baik lagi, ada beberapa saran yang mungkin dapat membantu
dalam mengembangkan optimasi beton bertulang pada struktur portal ruang
yaitu:
\begin{enumerate}[label=\arabic*.]
\item Dalam memilih bahasa pemrograman, sebaiknya digunakan bahasa yang
      dapat menembus memori lebih dari 64K, karena dalam pemrograman
      optimasi beton bertulang pada struktur portal ruang dibutuhkan
      dimensi sebesar mungkin, selain itu untuk sistem operasi dapat
      dicoba menggunakan sistem operasi yang sudah berbasiskan 64 bit
      atau yang \textit{flat memory}, misalnya menggunakan kernel Linux.
\item Metoda optimasi polihedron fleksibel perlu dikembangkan lagi,
      terutama dalam proses penelusurannya, sehingga metoda ini dapat
      menjadi lebih baik lagi.
\end{enumerate}


% Daftar Pustaka -- transcribed verbatim from Daftar/Daftar Pustaka.doc
\begin{thebibliography}{99}
\addcontentsline{toc}{chapter}{DAFTAR PUSTAKA}

\bibitem{balfour1986} Balfour, James A. D., 1986, \textit{Computer Analysis of Structural Frameworks}, Collins, London.

\bibitem{dipohusodo1994} Dipohusodo, I., 1994, \textit{Struktur Beton Bertulang}, Departemen Pekerjaan Umum RI, Gramedia Pustaka Utama, Jakarta.

\bibitem{harsoyo1976} Harsoyo, 1976, \textit{Seri Bina Bangunan, Konstruksi Beton Bertulang, Gedung Bertingkat II}, R. Sugihardjo BAE, Yogyakarta.

\bibitem{hulsemosley1989} Hulse, R. and Mosley, W.H., 1989, \textit{Reinforced Concrete Design by Computer}, Macmillan, Singapore.

\bibitem{kirsch1981} Kirsch, U., 1981, \textit{Optimum Structural Design}, McGraw-Hill Company.

\bibitem{laiachenbach1973} Lai, Y. S., and Achenbach, J. D., 1973, Direct Search Optimization Method, \textit{Journal of Structural Division, ASCE}, Vol. 98, pp. 119--131.

\bibitem{lev1977} Lev, O. E., 1977, A Structural Optimization Solution to a Branch and Bound Problem, \textit{Quarterly of Applied Mathematics}, Vol. 34, pp. 365--371.

\bibitem{leveyfu1979} Levey, G. C., and Fu, K. C., 1979, A Method in Discrete Frame Optimization and Its Outlook, \textit{International Journal of Computers and Structures}, Vol. 10, pp. 2177--2197.

\bibitem{macgregor1997} MacGregor, James G., 1997, International Edition, \textit{Reinforced Concrete, Mechanics and Design}, Third Edition, Prentice-Hall, New Jersey.

\bibitem{morris1982} Morris, A. J., 1982, \textit{Foundations of Structural Optimization: A Unified Approach}, John Wiley \& Sons, New York.

\bibitem{nawy1990} Nawy, E., 1990, \textit{Beton Bertulang Suatu Pendekatan Dasar}, Terjemahan, PT Eresco, Bandung.

\bibitem{sksni1991} SK SNI T-15-1991-03, 1991, \textit{Standar Tata Cara Perhitungan Struktur Beton untuk Bangunan Gedung}, Departemen Pekerjaan Umum, Yayasan Lembaga Penyelidikan Masalah Bangunan, Bandung.

\bibitem{sutedjomichael1997} Sutedjo, B. dan Michael, 1997, \textit{Algoritma dan Teknik Pemrograman}, Andi, Yogyakarta.

\bibitem{toakley1968} Toakley, A. R., 1968, Optimum Design Using Available Section, \textit{Journal of Structural Division, ASCE}, Vol. 34, pp. 1219--1241.

\bibitem{wangsalmon1990} Wang, C.K. dan Salmon C.G., 1990, \textit{Desain Beton Bertulang}, Edisi keempat, Terjemahan, Erlangga, Jakarta.

\bibitem{wahyudirahim1997} Wahyudi, L. dan Syahril A.R., 1997, \textit{Struktur Beton Bertulang Standar Baru SNI T-15-1991-03}, PT Gramedia Pustaka Utama, Jakarta.

\bibitem{weavergere1980} Weaver, W. Jr. dan James M. Gere, \textit{Analisa Matriks Untuk Struktur Rangka}, Edisi kedua, Terjemahan, Erlangga, Jakarta.

\bibitem{wibowo1996} Wibowo, F.N., 1996, \textit{Perkembangan Metoda Optimasi pada Teknik Sipil}, diktat kuliah Metoda Optimasi Struktur, Universitas Atma Jaya Yogyakarta, Yogyakarta.

\bibitem{wu1986} Wu, Peiming, 1986, \textit{Structural Optimization with Nonlinear and Discrete Programming}, University Microfilm International, Michigan.

\bibitem{viskusuma1993} Vis, W.C. dan Kusuma G.H., 1993, \textit{Dasar-dasar Perencanaan Beton Bertulang}, Erlangga, Jakarta.

\bibitem{young1989} Young, Warren C., 1989, \textit{Roark's Formulas for Stress \& Strain}, McGraw-Hill, New York.

\end{thebibliography}


% Lampiran (appendices): the 38 appendices in Teori/Lampiran/ are mostly
% verbatim source-code listings that already live, as real working code, in
% Optimasi Beton/Source/*.hpp -- they are referenced here rather than
% retranscribed. See LEGACY_TRANSCRIPTION_NOTES.md.

\end{document}
